\documentclass[aip,jcp,reprint,noshowkeys,superscriptaddress]{revtex4-1}
\usepackage{graphicx,dcolumn,bm,xcolor,microtype,multirow,amscd,amsmath,amssymb,amsfonts,physics,longtable,wrapfig}
\usepackage[version=4]{mhchem}

\usepackage[utf8]{inputenc}
\usepackage[T1]{fontenc}
\usepackage{txfonts}

\usepackage[
    colorlinks,
    linkcolor={red!50!black},
    citecolor={red!70!black},
    urlcolor={red!80!black}
	]{hyperref}
\urlstyle{same}

\newcommand{\ie}{\textit{i.e.}}
\newcommand{\eg}{\textit{e.g.}}
\newcommand{\alert}[1]{\textcolor{red}{#1}}
\usepackage[normalem]{ulem}
\newcommand{\titou}[1]{\textcolor{red}{#1}}
\newcommand{\trashPFL}[1]{\textcolor{\red}{\sout{#1}}}
\newcommand{\PFL}[1]{\titou{(\underline{\bf PFL}: #1)}}

\newcommand{\cre}[1]{\Hat{a}_{#1}^{\dag}}
\newcommand{\ani}[1]{\Hat{a}_{#1}}
\newcommand{\ca}[2]{\Hat{a}^{#1}_{#2}}

\newcommand{\mc}{\multicolumn}
\newcommand{\fnm}{\footnotemark}
\newcommand{\fnt}{\footnotetext}
\newcommand{\tabc}[1]{\multicolumn{1}{c}{#1}}
\newcommand{\QP}{\textsc{quantum package}}
\newcommand{\T}[1]{#1^{\intercal}}

% coordinates
\newcommand{\br}{\boldsymbol{r}}
\newcommand{\bx}{\boldsymbol{x}}
\newcommand{\bI}{\boldsymbol{1}}
\newcommand{\cK}{\mathcal{K}}
\newcommand{\bG}{\boldsymbol{G}}
\newcommand{\bH}{\boldsymbol{H}}
\newcommand{\bM}{\boldsymbol{M}}
\newcommand{\bE}{\boldsymbol{E}}
\newcommand{\bSigma}{\boldsymbol{\Sigma}}

% orbital energies
\newcommand{\eps}{\epsilon}
\newcommand{\irr}{\text{irr}}

% shortcuts for greek letters
\newcommand{\si}{\sigma}
\newcommand{\la}{\lambda}
\newcommand{\ii}{\mathrm{i}}
\newcommand{\up}{\uparrow}
\newcommand{\dw}{\downarrow}
\newcommand{\upup}{\up\up}
\newcommand{\updw}{\up\dw}
\newcommand{\dwup}{\dw\up}
\newcommand{\dwdw}{\dw\downarrow}

\newcommand{\h}{\text{1h}}
\newcommand{\p}{\text{1p}}
\newcommand{\hp}{\text{1h+1p}}
\newcommand{\hhp}{\text{2h1p}}
\newcommand{\pph}{\text{2p1h}}
\newcommand{\ph}{\text{ph}}
\newcommand{\eh}{\text{eh}}
\newcommand{\he}{\text{he}}
\newcommand{\ee}{\text{ee}}
\newcommand{\teh}{\overline{\text{eh}}}
\newcommand{\pp}{\text{pp}}
\newcommand{\hh}{\text{hh}}
\newcommand{\xc}{\text{xc}}


% addresses
\newcommand{\LCPQ}{Laboratoire de Chimie et Physique Quantiques (UMR 5626), Universit\'e de Toulouse, CNRS, UPS, France}


% statement-status markers used throughout the notes
\newcommand{\Exact}{\textbf{Exact identity.}\ }
\newcommand{\Established}{\textbf{Established result.}\ }
\newcommand{\Proposed}{\textbf{Proposed approximation.}\ }

\begin{document}

\title{Multireference Hedin Equations}

%\author{Pierre-Fran\c{c}ois \surname{Loos}}
%	\email{loos@irsamc.ups-tlse.fr}
%	\affiliation{\LCPQ}

\begin{abstract}
We develop a formal and practical framework for multireference Hedin equations
built around a CASSCF wave function and the Dyall Hamiltonian.  The first
objective is to formulate an exact Dyall-reference Hedin hierarchy and to show
that the multireference $GW$ approximation of Wang, Fang, and Li can be
understood as a channel-selective vertex and self-energy truncation of this
hierarchy.  The second objective is to go beyond that approximation at minimal
additional cost by restoring the mixed Hall self-energy blocks, screening them
with a static kernel consistent with the parent MR-$GW$ approximation, and
embedding the resulting couplings in a common Hermitian representation.  The
latter construction is causal, has a positive-semidefinite spectral function,
preserves the zeroth spectral moment, and restores the exact first-order
contribution to the first spectral moment.  We also give a matrix-free
computational formulation whose leading dense orbital scaling is unchanged
relative to the parent spectral MR-$GW$ implementation.
\end{abstract}

\maketitle


\section{Multireference Green's Function methods}

Assuming the exact self-energy, the one-particle Green's function can be obtained from the Dyson equation
\begin{equation}
\bG(\omega) = \bG_0(\omega) + \bG_0(\omega) \bSigma(\omega) \bG(\omega),
\end{equation}
where $\bG_0$ is the non-interacting Green's function and $\bSigma$ is the exact self-energy.
Now we can partition the self-energy into two contributions $\bSigma = \bSigma^{\text{init}} + \bSigma^{\text{corr}}$, where $\bSigma^{\text{init}}$ captures the contributions due to some initial treatment of correlation effects and $\bSigma^{\text{corr}}$ contains all the effects that are not contained within $\bSigma^{\text{init}}$.

The Dyson equation becomes 
\begin{align}
\bG(\omega) &= \bG_0(\omega) + \bG_0(\omega) \left( \bSigma^{\text{init}}(\omega) + \bSigma^{\text{corr}}(\omega) \right) \bG(\omega).
\end{align}
After inversion, one finds 
\begin{align}
\bG(\omega) = \frac{1}{\bG_0^{-1}(\omega) - \bSigma^{\text{init}}(\omega) - \bSigma^{\text{corr}}(\omega)},
\end{align}
where the new initial Green's function is defined as 
\begin{align}
  \bG_{\text{init}}(\omega) = \frac{1}{\bG_0^{-1}(\omega) - \bSigma^{\text{init}}(\omega)} .
\end{align}
The full Green's function can be expressed in terms of the initial Green's function and the correlation self-energy as
\begin{align}
\bG(\omega) &= \frac{1}{\bG_{\text{init}}^{-1}(\omega) - \bSigma^{\text{corr}}(\omega)} \nonumber \\
&= \bG_{\text{init}}(\omega) + \bG_{\text{init}}(\omega) \bSigma^{\text{corr}}(\omega) \bG(\omega).
\end{align}

Next we assume that the initial $\bSigma^{\text{init}}$ is obtained from an active space calculation, e.g., for a set of orbitals $\{ \phi_p \}$ defining the active space.
The total orbital space can be partitioned into active and inactive orbitals, so that the ground-state wavefunction is 
\begin{align}
  | \Psi_0 \rangle = | \Phi_\mathrm{i} \rangle \otimes | \Psi_\mathrm{a} \rangle,
\end{align}
where $\mathrm{i}$ and $\mathrm{a}$ denote the inactive and active orbitals, respectively.

Within the chosen partitioning, the initial Green's function $\bG_{\text{init}}$ can also be directly determined without resorting to the Dyson equation, for example, by using the spectral representation of the active space Hamiltonian in the $N\pm1$ particle sectors in combination with the coupling between the inactive orbitals to the active space. 
$\bG_{\text{init}}$ reads
\begin{align}
\bG_{\text{init}}(\omega) = 
\begin{pmatrix}
  \mathbf{1} & 0 \\
  0 & \bM
\end{pmatrix}
\begin{pmatrix}
  \omega \mathbf{1} - \mathbf{F}^\mathrm{HF}_\mathrm{i,i} & \bH_\mathrm{i, a}\\
   \bH^\dagger_\mathrm{i, a} & \omega \mathbf{1} - \bE^{N \pm 1}
\end{pmatrix}^{-1}
\begin{pmatrix}
  \mathbf{1} & 0 \\
  0 & \bM^\dagger
\end{pmatrix},
\label{eq:Ginit}
\end{align}
with $\bM$ denoting the Dyson amplitudes within the active space, e.g.,
\begin{align}
  M^{N+1}_{p,\mu} &= \left( p | R^{N+1}_\mu \right) = \langle \Psi_0 | \left[ a_p, \hat{R}^{N+1}_\mu \right] | \Psi_0 \rangle = \langle \Psi_0 | a_p | \Psi_\mu^{N+1} \rangle, \nonumber \\
  M^{N-1}_{p,\mu} &= \left( p | R^{N-1}_\mu \right) = \langle \Psi_0 | \left[ a_p^\dagger, \hat{R}^{N-1}_\mu \right] | \Psi_0 \rangle = \langle \Psi_0 | a_p^\dagger | \Psi_\mu^{N-1} \rangle,  \nonumber
\end{align}
where $\left( a | b \right)$ denotes the binary product 
\begin{align}
  \left( a | b \right) = \langle \Psi_0 | \left[ a^\dagger,  b \right] | \Psi_0 \rangle.
\end{align}
The excitation operators $\hat{R}^{N\pm1}_\mu$ act only within the active space, and satisfies the `killer condition', i.e.,
\begin{align}
  \langle \Psi_0 | \hat{R}^{N+1}_\mu  = 0, \quad  \langle \Psi_0 | \hat{R}^{N-1}_\mu = 0.
\end{align}
Finally, $\bE^{N \pm 1}$ denotes a diagonal matrix with the excitation energies of the active space in the $N\pm1$ particle sectors on the diagonal.

Note that there are no Dyson amplitudes that couple the active and inactive spaces directly.
However, there are Hamiltonian elements $ \bH_\mathrm{i, a}$ that couple the inactive and active spaces.
Analytic expressions can be obtained by defining the excitation manifolds of the inactive orbitals as
\begin{align}
  \left\{ \hat{C}^{\dagger} \right\} = \left\{ a_i, a_a  \right\}.
  \label{eq:C1}
\end{align}
The matrix elements $ \bH_\mathrm{i, a}$  read 
\begin{align}
  \bH_\mathrm{i, a}  = \langle \Psi_0  | \left[ \hat{C}, \left[ \hat{H}, \hat{R}^{N \pm 1} \right] \right] | \Psi_0 \rangle.
\end{align}
\textcolor{red}{\textbf{One has to be a little bit careful. There might be some problem with the sign convention}}
Since the Hamiltonian contains only one- and two-body terms, only a limited number of excitation manifolds in the excitation operators $\hat{R}^{N \pm 1}$ need to be explicitly considered.

With this, one can rewrite Eq.~\eqref{eq:Ginit} as
\begin{align}
  \bG_{\text{init}}(\omega) &= \nonumber \\
 &\begin{pmatrix}
    \bG_\mathrm{i}(\omega) & \bG_\mathrm{i}(\omega) \bSigma_\mathrm{i,a}(\omega) \\
    \bSigma_\mathrm{a,i}(\omega) \bG_\mathrm{i}(\omega) & \bG_\mathrm{A} + \bSigma_\mathrm{a,i}(\omega) \bG_\mathrm{i}(\omega) \bSigma_\mathrm{i,a}(\omega) \\
  \end{pmatrix},
  \label{eq:Ginit2}
\end{align}
where 
\begin{align}
  \bSigma_\mathrm{i,a}(\omega) = - \bH_\mathrm{i, a} \left( \omega - \bE^{N \pm 1} \right)^{-1} \bM \nonumber ,
\end{align}
and
\begin{align}
  \bG_\mathrm{i}(\omega) = \left[ \bG^{-1}_{\mathrm{i,0}}(\omega) - \bSigma_\mathrm{i,i}(\omega) \right]^{-1} = \bG_{\mathrm{i,0}}(\omega) + \tilde{\bSigma}_\mathrm{i,i}(\omega),
\end{align}
with 
\begin{align}
    \bSigma_\mathrm{i,i}(\omega) = \bH_\mathrm{i, a} \left( \omega - \bE^{N \pm 1} \right)^{-1} \bH_\mathrm{a, i}.
\end{align}
Eq.~\eqref{eq:Ginit2} can be compactly written as
\begin{align}
  \bG_\mathrm{init} = \left( \mathbf{1} - \bSigma_\mathrm{a,i}(\omega) \right)^{-1} \left( \bG_{\mathrm{i,0}}(\omega) + \tilde{\bSigma}_\mathrm{i,i}(\omega)  \right) \left( \mathbf{1} - \bSigma_\mathrm{i,a}(\omega) \right)^{-1},
  \label{eq:Ginit3}
\end{align}
with
\begin{align}
   \left( \mathbf{1} - \bSigma_\mathrm{a,i}(\omega) \right)^{-1} = 
   \begin{pmatrix}
    \mathbf{1} & 0  \\
    - \bSigma_\mathrm{a,i}(\omega) & \mathbf{1}
   \end{pmatrix}^{-1}
   = 
  \begin{pmatrix}
    \mathbf{1} & 0  \\
    \bSigma_\mathrm{a,i}(\omega) & \mathbf{1}
   \end{pmatrix}.
\end{align}

Eq.~\eqref{eq:Ginit3} has the same structure as the initial Dyson equation in the MR-$GW$ paper. 
However, in this derivation the meaning of the different terms becomes clearer.
In particular, we have contributions reflecting the fact that an initial correlated state $| \Psi_0 \rangle$ is used.
Furthermore, all remaining contributions in $\bSigma^\mathrm{corr.}$ describe two additional correlation effects missing in $G_{\mathrm{init}}$:
\begin{itemize}
\item Effects within the inactive space itself.
\item Couplings between the inactive and active spaces that go beyond the operator manifold defined in Eq.~\eqref{eq:C1}.
\end{itemize}

\clearpage
\newpage

% ============================================================
% Task 1: Exact Dyall-reference Hedin hierarchy
% ============================================================

\PFL{These notes have been automatically generated by ChatGPT after a very lengthy convservation.
They need to be carefully understood, checked, and polished.}
%%%%%%%%%%%%%%%%%%%%%%%%%
\section{Introduction and objectives}
\label{sec:introduction}
%%%%%%%%%%%%%%%%%%%%%%%%%

The aim of these notes is to develop a multireference extension of Hedin's
framework in which the zeroth-order Hamiltonian is the Dyall Hamiltonian
associated with a CASSCF wave function.  The central difficulty is that the
Dyall reference is interacting and therefore non-Gaussian: its exact
generating functional contains connected cumulants $C_n^D$ of all orders.
Consequently, the standard five-equation Hedin construction cannot simply be
reused without modification.

Two complementary objectives are pursued.

\paragraph{Objective I: formal embedding of existing MR-$GW$.}
We first construct an exact Dyall-reference Hedin hierarchy by combining the
exact correlated Dyall generating functional with an enlarged source-space
representation inspired by Brouder's treatment of general correlated
references \cite{Brouder2007}.  We then verify its first-order expansion
against the four generalized self-energy blocks derived by Wang, Fang, and Li
\cite{WangFangLi2026}, and show that their practical MR-$GW$ scheme can be
understood as a channel-selective truncation of the resulting hierarchy.

\paragraph{Objective II: a minimal nonzero-block extension.}
We next restore the lowest-cost information discarded by practical MR-$GW$,
namely the mixed Hall blocks $\Sigma^{12}$ and $\Sigma^{21}$.  We construct a
static screened kernel consistent with the parent MR-$GW$ self-energy under a
frozen-$W$ approximation, and embed the resulting mixed couplings together
with the dynamic MR-$GW$ self-energy in a common Hermitian pole-space
representation.  This defines a causal positive-semidefinite completion that
preserves the zeroth spectral moment and is exact through first residual order
for the first moment.

Throughout the notes, the Hamiltonian partition
\begin{equation}
  \hat H=\hat H_D+\hat V_R
  \label{eq:global-partition}
\end{equation}
is used without exception.  The active-space interaction is treated
nonperturbatively in $\hat H_D$, while only the residual interaction
$\hat V_R$ is screened and expanded.
%%%%%%%%%%%%%%%%%%%%%%%%%
\subsection{Status of statements}
%%%%%%%%%%%%%%%%%%%%%%%%%

To keep the logical status of the various equations explicit, we use the
following terminology.

\Exact denotes an identity within the exact Dyall-reference or enlarged-space
formalism.

\Established denotes a result already present in the literature, most notably
the generalized Dyson construction of Hall and Brouder and the first-order
self-energy/MR-$GW$ formulas of Wang, Fang, and Li.

\Proposed denotes an approximation introduced in these notes.  In particular,
the static mixed kernel and its finite Hermitian/PSD resummation belong to
this category.  Some algebraic properties of a proposed construction, such as
its Hermiticity or its spectral-moment consequences, may of course be exact
once that construction has been defined.
%%%%%%%%%%%%%%%%%%%%%%%%%
\subsection{Organization}
%%%%%%%%%%%%%%%%%%%%%%%%%

Sections~\ref{sec:objective-I} and \ref{sec:objective-II} correspond to the two
objectives above.  A consolidated set of limiting cases is collected in
Sec.~\ref{sec:consistency}, followed by conclusions and the main directions
for further work.
%%%%%%%%%%%%%%%%%%%%%%%%%
\section{Objective I: Dyall-reference Hedin theory and the existing MR-$GW$ scheme}
\label{sec:objective-I}
%%%%%%%%%%%%%%%%%%%%%%%%%
\subsection{Exact Dyall-reference Hedin hierarchy}
\label{sec:dyall-hedin}
%%%%%%%%%%%%%%%%%%%%%%%%%

\Exact The equations in this subsection are identities of the enlarged
Dyall-reference theory once the auxiliary cumulant interaction is defined by
the source-transfer identity introduced below.  Brouder's construction for a
general correlated state and its connection to Hall's generalized Dyson
equation are established in the literature; applying the same algebraic
source transfer to the exact generating functional of the interacting Dyall
reference is the extension developed here.


Our starting point is a partition of the physical Hamiltonian into a
Dyall zeroth-order Hamiltonian and a residual interaction,
\begin{equation}
  \hat H = \hat H_D + \hat V_R
  \label{eq:dyall-partition}
\end{equation}
The Dyall zeroth-order Hamiltonian is associated with a chosen CASSCF
reference and contains the full interaction within the active orbital space.
In contrast with the usual Hartree--Fock reference used in conventional
many-body perturbation theory, $\hat H_D$ is therefore itself an interacting
Hamiltonian.  We use ``Dyall Hamiltonian'' in this sense throughout, rather
than identifying it with the conventional CASPT2 zeroth-order Hamiltonian.

This distinction is essential.  The exact generating functional
associated with $\hat H_D$ is non-Gaussian and contains connected
$n$-particle Green's functions of arbitrary order.  The purpose of this
section is to show how this non-Gaussian reference can nevertheless be
embedded into a Hedin-like framework by introducing an enlarged source
space following the general strategy introduced by Brouder
\cite{Brouder2007}.  The physical $11$ sector of the resulting theory
reproduces Hall's generalized Dyson equation and therefore provides the
natural formal setting for the multireference Green's-function theory of
Wang, Fang, and Li \cite{WangFangLi2026}.


\subsubsection{Non-Gaussian Dyall reference}

Let
\begin{equation}
  G_D(1,2)
  =
  -i
  \braket{
    T_C
    \hat\psi_D(1)
    \hat\psi_D^\dagger(2)
  }
\end{equation}
denote the exact one-particle Green's function of the Dyall Hamiltonian,
where $C$ denotes the chosen time contour.

Introducing Grassmann sources $\bar\eta$ and $\eta$, the logarithm of the
exact Dyall generating functional can be decomposed as
\begin{equation}
  \ln Z_D[\bar\eta,\eta]
  =
  -i \bar\eta G_D \eta
  +
  \rho_D[\bar\eta,\eta]
  \label{eq:dyall-generating-functional}
\end{equation}
The first term contains the exact one-particle propagator, while
$\rho_D$ contains all connected reference correlations beyond the
one-particle level.

Particle-number conservation implies the expansion
\begin{equation}
  \rho_D
  =
  \rho_D^{(2)}
  +
  \rho_D^{(3)}
  +
  \cdots
\end{equation}
where $\rho_D^{(n)}$ is homogeneous of degree $n$ in $\bar\eta$ and
degree $n$ in $\eta$.  Its coefficients are the connected Dyall
$n$-particle Green's functions,
\begin{equation}
  C_n^D
  \equiv
  G_{D,c}^{(n)}
\end{equation}

For a Gaussian reference,
\begin{equation}
  C_{n\ge 2}^D = 0
\end{equation}
whereas for a genuine CASSCF/Dyall reference one generally has
\begin{equation}
  C_2^D
  C_3^D
  \ldots
  \neq 0
\end{equation}
This is the precise field-theoretical meaning of the non-Gaussian
character of the multireference zeroth-order problem.


\subsubsection{Enlarged Gaussian representation}

Brouder showed that the cumulants of a general correlated state can be
transferred from the reference generating functional to auxiliary
interaction vertices by introducing an enlarged source space
\cite{Brouder2007}.  In Brouder's original formulation the underlying
forward/backward contour (or source) structure is explicit; below we suppress
these contour indices and retain only the compact $2\times2$ auxiliary-space
structure.  We use the same source-doubling idea algebraically, but apply it
to the exact generating functional of the interacting Dyall reference rather
than to a quadratic reference Hamiltonian.  The following representation is
therefore exact once the auxiliary cumulant interaction is defined by
Eq.~\eqref{eq:cumulant-transfer}.

Introduce a second pair of Grassmann sources,
\begin{equation}
  \bar\xi
  \qquad
  \xi
\end{equation}
and define
\begin{equation}
  \overline{\mathcal J}
  =
  \begin{pmatrix}
    \bar\eta & \bar\xi
  \end{pmatrix}
  \qquad
  \mathcal J
  =
  \begin{pmatrix}
    \eta \\
    \xi
  \end{pmatrix}
\end{equation}
The enlarged Gaussian reference propagator is chosen as
\begin{equation}
  \mathbb G_D
  =
  \begin{pmatrix}
    G_D & -i I \\
    -i I & 0
  \end{pmatrix}
  \label{eq:super-dyall-gf}
\end{equation}
with inverse
\begin{equation}
  \mathbb G_D^{-1}
  =
  \begin{pmatrix}
    0 & i I \\
    i I & G_D
  \end{pmatrix}
  \label{eq:super-dyall-gf-inverse}
\end{equation}

The auxiliary cumulant interaction $\widehat{\mathcal C}_D$ is defined
by the exact source-transfer identity
\begin{equation}
  e^{-i\bar\eta G_D\eta + \rho_D[\bar\eta,\eta]}
  =
  \left.
  e^{-i\widehat{\mathcal C}_D}
  \exp
  \left[
    -i
    \overline{\mathcal J}
    \mathbb G_D
    \mathcal J
  \right]
  \right|_{\xi=\bar\xi=0}
  \label{eq:cumulant-transfer}
\end{equation}
Equation~\eqref{eq:cumulant-transfer} may be regarded as the definition
of $\widehat{\mathcal C}_D$.  It fixes all Grassmann signs and
combinatorial factors unambiguously.

Expanding $\widehat{\mathcal C}_D$ generates a hierarchy of auxiliary
interaction vertices carrying
\begin{equation}
  C_2^D
  \quad
  C_3^D
  \quad
  \ldots
\end{equation}
The non-Gaussianity of the Dyall reference has therefore been moved from
the quadratic part of the theory into a set of many-leg auxiliary
vertices.


\subsubsection{Enlarged Dyson equation}

The residual physical interaction acts in the physical sector, whereas
the reference cumulant interaction acts in the auxiliary sector.  The
enlarged theory may therefore be written schematically as
\begin{equation}
  \mathcal V_{\mathrm{ext}}
  =
  V_R[\bar\psi_1,\psi_1]
  +
  \mathcal C_D[\bar\psi_2,\psi_2]
  \label{eq:extended-interaction}
\end{equation}

Let $\mathbb G$ denote the exact enlarged Green's function.  Since the
quadratic part of the enlarged theory is Gaussian, one obtains an
ordinary Dyson equation,
\begin{equation}
  \mathbb G^{-1}
  =
  \mathbb G_D^{-1}
  -
  \overline{\boldsymbol\Sigma}
  \label{eq:super-dyson}
\end{equation}

Writing
\begin{equation}
  \overline{\boldsymbol\Sigma}
  =
  \begin{pmatrix}
    \overline\Sigma_{11} &
    \overline\Sigma_{12}
    \\
    \overline\Sigma_{21} &
    \overline\Sigma_{22}
  \end{pmatrix}
\end{equation}
the physical Green's function is simply
\begin{equation}
  G = \mathbb G_{11}
\end{equation}

The correspondence with Hall's generalized self-energy blocks is
\begin{subequations}
  \label{eq:hall-mapping}
  \begin{align}
    \Sigma^{11}
    &=
    \overline\Sigma_{11}
    \\
    \Sigma^{12}
    &=
    -i \overline\Sigma_{12}
    \\
    \Sigma^{21}
    &=
    -i \overline\Sigma_{21}
    \\
    \Sigma^{22}
    &=
    -\overline\Sigma_{22}
  \end{align}
\end{subequations}
Block inversion of Eq.~\eqref{eq:super-dyson} yields
\begin{equation}
  M
  =
  (1-\Sigma^{21})^{-1}
  \left(
    G_D+\Sigma^{22}
  \right)
  (1-\Sigma^{12})^{-1}
  \label{eq:hall-M}
\end{equation}
followed by
\begin{equation}
  G
  =
  M
  +
  M\Sigma^{11}G
  \label{eq:hall-dyson}
\end{equation}
Equations~\eqref{eq:hall-M} and \eqref{eq:hall-dyson} are precisely
Hall's generalized Dyson structure and coincide with the formulation
used in Ref.~\cite{WangFangLi2026}.


\subsubsection{Physical source and enlarged vertex}

Introduce a bilocal physical one-body source
\begin{equation}
  \hat U
  =
  \int d(12)\
  U(1,2)
  \hat\psi^\dagger(1)
  \hat\psi(2)
\end{equation}
The source acts only in the physical sector, so that
\begin{equation}
  \mathbb U
  =
  PUP
  \qquad
  P
  =
  \begin{pmatrix}
    I & 0 \\
    0 & 0
  \end{pmatrix}
\end{equation}

We choose the source convention
\begin{equation}
  \mathbb G^{-1}[U]
  =
  \mathbb G_D^{-1}
  -
  \mathbb U
  -
  \overline{\boldsymbol\Sigma}[U]
  \label{eq:source-dyson}
\end{equation}

The exact enlarged three-point vertex is defined as
\begin{equation}
  \overline{\boldsymbol\Gamma}
  =
  -
  \frac{
    \delta\mathbb G^{-1}
  }{
    \delta U
  }
  \label{eq:super-vertex-def}
\end{equation}
Using
\begin{equation}
  \delta\mathbb G
  =
  -
  \mathbb G
  \left(
    \delta\mathbb G^{-1}
  \right)
  \mathbb G
\end{equation}
one obtains the exact Schwinger relation
\begin{equation}
  \frac{
    \delta\mathbb G
  }{
    \delta U
  }
  =
  \mathbb G
  \overline{\boldsymbol\Gamma}
  \mathbb G
  \label{eq:super-schwinger}
\end{equation}


\subsubsection{Superspace Bethe--Salpeter equation}

From Eq.~\eqref{eq:source-dyson},
\begin{equation}
  \overline{\boldsymbol\Gamma}
  =
  P
  +
  \frac{
    \delta\overline{\boldsymbol\Sigma}
  }{
    \delta U
  }
\end{equation}
Introduce the exact enlarged irreducible kernel
\begin{equation}
  \overline K_{AB;CD}
  =
  \frac{
    \delta\overline\Sigma_{AB}
  }{
    \delta\mathbb G_{CD}
  }
  \label{eq:super-kernel}
\end{equation}
The chain rule then gives
\begin{equation}
  \overline{\boldsymbol\Gamma}
  =
  P
  +
  \overline K
  \
  \mathbb G
  \overline{\boldsymbol\Gamma}
  \mathbb G
  \label{eq:super-bse}
\end{equation}
Equation~\eqref{eq:super-bse} is the exact Bethe--Salpeter equation in
the enlarged source space.


\subsubsection{Exact Dyall Schwinger identity}

The reference response can also be obtained directly, without invoking
the enlarged formalism.  Differentiating the exact Dyall Green's
function with respect to the physical source gives
\begin{equation}
  \begin{aligned}
  \frac{
    \delta G^D_{pq}(t,t')
  }{
    \delta U_{rs}(\tau)
  }
  ={}&
  G^D_{pr}(t,\tau)
  G^D_{sq}(\tau,t')
  \\
  &-
  [C_2^D]_{sp,rq}
  (\tau,t,\tau^+,t')
  \end{aligned}
  \label{eq:dyall-schwinger}
\end{equation}
In compact notation,
\begin{equation}
  \frac{\delta G_D}{\delta U}
  =
  G_DG_D
  -
  C_2^{D,ph}
  \label{eq:dyall-schwinger-compact}
\end{equation}
The same structure appears explicitly in the first-order expansion of
Ref.~\cite{WangFangLi2026}.

Consequently, the enlarged reference vertex must satisfy
\begin{equation}
  \left[
    \mathbb G_D
    \overline{\boldsymbol\Gamma}_D
    \mathbb G_D
  \right]_{11}
  =
  G_DG_D
  -
  C_2^{D,ph}
  \label{eq:reference-vertex-normalization}
\end{equation}
This relation fixes the full physical projection of the reference
vertex.  It does \emph{not} imply that any individual enlarged-space
component, such as $\Gamma_{22,D}$, can be identified directly with
$C_2^D$.


\subsubsection{Effective physical Dyall vertex}

It is convenient to define an effective physical reference vertex
$\Gamma_D$ through
\begin{equation}
  G_D
  \Gamma_D
  G_D
  \equiv
  \left[
    \mathbb G_D
    \overline{\boldsymbol\Gamma}_D
    \mathbb G_D
  \right]_{11}
\end{equation}
Equation~\eqref{eq:reference-vertex-normalization} then gives
\begin{equation}
  G_D
  \Gamma_D
  G_D
  =
  G_DG_D
  -
  C_2^{D,ph}
  \label{eq:physical-dyall-vertex}
\end{equation}
Formally,
\begin{equation}
  \Gamma_D
  =
  1
  -
  G_D^{-1}
  C_2^{D,ph}
  G_D^{-1}
  \label{eq:dyall-vertex}
\end{equation}
where all integrations and particle--hole index permutations are
understood.

For a correlated reference,
\begin{equation}
  \Gamma_D\neq1
\end{equation}
In the Gaussian limit,
\begin{equation}
  C_2^D\rightarrow0
\end{equation}
and therefore
\begin{equation}
  \Gamma_D\rightarrow1
\end{equation}


\subsubsection{Polarization and screened residual interaction}

The exact physical irreducible polarization is defined by
\begin{equation}
  P_R
  =
  -i
  \left[
    \mathbb G
    \overline{\boldsymbol\Gamma}
    \mathbb G
  \right]_{11}^{ph}
  \label{eq:mr-polarization}
\end{equation}

At the Dyall-reference level,
\begin{equation}
  P_D
  =
  -i
  \left(
    G_DG_D
    -
    C_2^{D,ph}
  \right)
  \label{eq:dyall-polarization}
\end{equation}
or
\begin{equation}
  P_D
  =
  -iG_DG_D
  +
  iC_2^{D,ph}
\end{equation}

Only the residual interaction should be screened because the
active--active interaction is already contained exactly in $H_D$.  We
therefore define
\begin{equation}
  W_R
  =
  v_R
  +
  v_RP_RW_R
  \label{eq:residual-screening}
\end{equation}

At the fixed Dyall-reference level,
\begin{equation}
  P_R\rightarrow P_D
\end{equation}
which gives
\begin{equation}
  W_D
  =
  v_R
  +
  v_RP_DW_D
  \label{eq:dyall-screening}
\end{equation}
This is the MR-RPA screened interaction employed in
Ref.~\cite{WangFangLi2026}.


\subsubsection{Exact Schwinger--Dyson hierarchy}

The interaction in the enlarged theory contains two qualitatively
different contributions.  Define
\begin{equation}
  \mathcal F_A
  =
  \frac{
    \delta \mathcal V_{\mathrm{ext}}
  }{
    \delta\bar\Psi_A
  }
\end{equation}
For the physical sector,
\begin{equation}
  \mathcal F_1
  =
  \frac{
    \delta V_R
  }{
    \delta\bar\psi_1
  }
\end{equation}
whereas in the auxiliary sector
\begin{equation}
  \mathcal F_2
  =
  \frac{
    \delta\mathcal C_D
  }{
    \delta\bar\psi_2
  }
\end{equation}

Introduce the interaction insertion
\begin{equation}
  \mathcal X_{AB}(1,2)
  =
  -i
  \braket{
    T_C
    \mathcal F_A(1)
    \Psi_B^\dagger(2)
  }
\end{equation}
The equation of motion gives
\begin{equation}
  \mathcal X
  =
  \overline{\boldsymbol\Sigma}
  \mathbb G
  \label{eq:super-sd-master}
\end{equation}

The physical row therefore satisfies
\begin{equation}
  \left(
    \overline{\boldsymbol\Sigma}
    \mathbb G
  \right)_{1B}
  =
  -i
  \braket{
    T_C
    \frac{
      \delta V_R
    }{
      \delta\bar\psi_1
    }
    \Psi_B^\dagger
  }
  \label{eq:physical-sd-row}
\end{equation}
while the auxiliary row satisfies
\begin{equation}
  \left(
    \overline{\boldsymbol\Sigma}
    \mathbb G
  \right)_{2B}
  =
  -i
  \braket{
    T_C
    \frac{
      \delta\mathcal C_D
    }{
      \delta\bar\psi_2
    }
    \Psi_B^\dagger
  }
  \label{eq:auxiliary-sd-row}
\end{equation}

Since
\begin{equation}
  \mathcal C_D
  =
  \mathcal C_D^{(2)}
  +
  \mathcal C_D^{(3)}
  +
  \cdots
\end{equation}
the auxiliary row contains the complete hierarchy
\begin{equation}
  C_2^D
  \quad
  C_3^D
  \quad
  C_4^D
  \quad
  \ldots
\end{equation}
For a term of the schematic form
\begin{equation}
  \mathcal C_D^{(n)}
  \propto
  \frac{1}{n!n!}
  C_n^D
  \bar\psi_2^n
  \psi_2^n
\end{equation}
functional differentiation gives
\begin{equation}
  \frac{
    \delta\mathcal C_D^{(n)}
  }{
    \delta\bar\psi_2
  }
  \propto
  \frac{1}{(n-1)!n!}
  C_n^D
  \bar\psi_2^{n-1}
  \psi_2^n
\end{equation}

Hence the exact multireference theory is not a closed set of five Hedin
equations.  Instead,
\begin{equation}
  \text{Dyall-Hedin theory}
  =
  \text{Hedin structure}
  +
  \text{reference-cumulant hierarchy}
  \label{eq:mr-hedin-summary}
\end{equation}


\subsubsection{Compact form of the Dyall-Hedin hierarchy}

The exact formal structure may be summarized as
\begin{subequations}
  \label{eq:dyall-hedin-set}
  \begin{align}
    \mathbb G^{-1}
    &=
    \mathbb G_D^{-1}
    -
    \mathbb U
    -
    \overline{\boldsymbol\Sigma}
    \\
    \frac{\delta\mathbb G}{\delta U}
    &=
    \mathbb G
    \overline{\boldsymbol\Gamma}
    \mathbb G
    \\
    \overline{\boldsymbol\Gamma}
    &=
    P
    +
    \overline K
    \mathbb G
    \overline{\boldsymbol\Gamma}
    \mathbb G
    \\
    P_R
    &=
    -i
    \left[
      \mathbb G
      \overline{\boldsymbol\Gamma}
      \mathbb G
    \right]_{11}^{ph}
    \\
    W_R
    &=
    v_R
    +
    v_RP_RW_R
    \\
    \left(
      \overline{\boldsymbol\Sigma}
      \mathbb G
    \right)_{1B}
    &=
    -i
    \braket{
      T_C
      \frac{
        \delta V_R
      }{
        \delta\bar\psi_1
      }
      \Psi_B^\dagger
    }
    \\
    \left(
      \overline{\boldsymbol\Sigma}
      \mathbb G
    \right)_{2B}
    &=
    -i
    \braket{
      T_C
      \frac{
        \delta\mathcal C_D
      }{
        \delta\bar\psi_2
      }
      \Psi_B^\dagger
    }
  \end{align}
\end{subequations}
The reference vertex is constrained by
\begin{equation}
  \left[
    \mathbb G_D
    \overline{\boldsymbol\Gamma}_D
    \mathbb G_D
  \right]_{11}
  =
  G_DG_D
  -
  C_2^{D,ph}
  \label{eq:dyall-reference-condition}
\end{equation}
%%%%%%%%%%%%%%%%%%%%%%%%%
\subsection{First-order recovery of the generalized self-energy blocks}
\label{sec:first-order-blocks}
%%%%%%%%%%%%%%%%%%%%%%%%%

\Established The first-order blocks derived below are the
Wang--Fang--Li results.  Their role here is to provide a component-level check
of the exact hierarchy and to fix all signs, time orderings, and numerical
prefactors.


We now verify explicitly that the Dyall-reference formalism developed
in Sec.~\ref{sec:dyall-hedin} reproduces the four first-order
generalized self-energy blocks derived by Wang, Fang, and Li
\cite{WangFangLi2026}.

Throughout this section,
\begin{equation}
  G_D \equiv G_0
\end{equation}
denotes the exact Green's function of the zeroth-order Hamiltonian,
and
\begin{equation}
  C_2^D \equiv G_0^c
  \qquad
  C_3^D \equiv G_0^{c,(3)}
\end{equation}
denote the connected two- and three-body zeroth-order Green's
functions.

We introduce a bookkeeping parameter $\lambda$,
\begin{equation}
  \hat H(\lambda)
  =
  \hat H_D
  +
  \lambda \hat V_R
\end{equation}
and expand
\begin{equation}
  G
  =
  G_D
  +
  \lambda G_1
  +
  \mathcal O(\lambda^2)
\end{equation}
together with
\begin{equation}
  \Sigma^{ij}
  =
  \lambda\Sigma_1^{ij}
  +
  \mathcal O(\lambda^2)
\end{equation}


\subsubsection{First-order expansion of Hall's generalized Dyson equation}

The generalized Dyson equation is
\begin{align}
  M
  &=
  (1-\Sigma^{21})^{-1}
  (G_D+\Sigma^{22})
  (1-\Sigma^{12})^{-1}
  \\
  G
  &=
  M
  +
  M\Sigma^{11}G
\end{align}

To first order,
\begin{equation}
  (1-\Sigma^{21})^{-1}
  =
  1+\lambda\Sigma_1^{21}
  +
  \mathcal O(\lambda^2)
\end{equation}
and similarly for $\Sigma^{12}$.  Hence
\begin{equation}
  M
  =
  G_D
  +
  \lambda
  \left(
    \Sigma_1^{21}G_D
    +
    G_D\Sigma_1^{12}
    +
    \Sigma_1^{22}
  \right)
  +
  \mathcal O(\lambda^2)
\end{equation}

Substitution into the second generalized Dyson equation gives
\begin{equation}
  G_1
  =
  G_D\Sigma_1^{11}G_D
  +
  G_D\Sigma_1^{12}
  +
  \Sigma_1^{21}G_D
  +
  \Sigma_1^{22}
  \label{eq:first-order-generalized-dyson}
\end{equation}

Equation~\eqref{eq:first-order-generalized-dyson} provides the
decomposition against which the direct perturbative expansion will be
matched.


\subsubsection{Residual perturbation}

The residual interaction is written as
\begin{equation}
  \hat V_R
  =
  \hat V_{1}
  +
  \hat V_{2}
\end{equation}
where
\begin{equation}
  \hat V_{1}
  =
  u_{pq}
  \hat p^\dagger
  \hat q
\end{equation}
and
\begin{equation}
  \hat V_{2}
  =
  \frac{1}{2}
  v_{pr,qs}
  \hat p^\dagger
  \hat q^\dagger
  \hat s
  \hat r
\end{equation}

Introducing the antisymmetrized interaction
\begin{equation}
  \bar v_{pr,qs}
  =
  v_{pr,qs}
  -
  v_{ps,qr}
\end{equation}
the two-body contribution may equivalently be evaluated with a
$1/4$ prefactor and the antisymmetrized interaction.


\subsubsection{First-order Gell-Mann--Low expansion}

The exact Green's function can be expanded as
\begin{equation}
  G_{pq}(t,t')
  =
  \frac{
    (-i)
    \braket{
      \mathcal T
      e^{-i\int dt_1\,\hat V_R(t_1)}
      \hat p(t)\hat q^\dagger(t')
    }_D
  }{
    \braket{
      \mathcal T
      e^{-i\int dt_1\,\hat V_R(t_1)}
    }_D
  }
\end{equation}

At first order,
\begin{equation}
  G_1
  =
  G_{1,1e}
  +
  G_{1,2e}
\end{equation}
where the subscripts refer to the one- and two-electron pieces of
$\hat V_R$.


\subsubsection{One-body perturbation}

The contribution from $\hat V_1$ is particularly transparent in the
present formulation.  Since $u_{rs}$ couples to a one-body operator,
the exact Dyall Schwinger identity derived previously immediately gives
\begin{equation}
  [G_{1,1e}]_{pq}(t,t')
  =
  u_{rs}
  \int dt_1\
  \frac{
    \delta[G_D]_{pq}(t,t')
  }{
    \delta U_{rs}(t_1)
  }
\end{equation}
Using
\begin{equation}
  \frac{\delta G_D}{\delta U}
  =
  G_DG_D-C_2^D
\end{equation}
one obtains
\begin{equation}
  \begin{aligned}
  [G_{1,1e}]_{pq}(t,t')
  =
  u_{rs}
  \int dt_1
  \Big[
  &
  [G_D]_{pr}(t,t_1)
  [G_D]_{sq}(t_1,t')
  \\
  &
  -
  [C_2^D]_{sp,rq}
  (t_1,t,t_1^+,t')
  \Big]
  \end{aligned}
  \label{eq:G1-onebody}
\end{equation}

The first term in Eq.~\eqref{eq:G1-onebody} contributes to
$\Sigma_1^{11}$, whereas the second contributes to
$\Sigma_1^{22}$.


\subsubsection{Two-body perturbation}

For the two-body perturbation, direct expansion followed by the
cumulant decomposition of the three-body time-ordered Green's function
gives
\begin{equation}
  \begin{aligned}
  [G_{1,2e}]_{pq}(t,t')
  =
  \frac{i}{4}
  \bar v_{rt,su}
  \int dt_1\
  \Big\{
  &
  [C_3^D]_{utp,qsr}
  (t_1^+,t_1,t,t',t_1^{++},t_1^{+++})
  \\[1mm]
  &
  -2
  [G_D]_{pr}(t,t_1)
  [C_2^D]_{ut,sq}
  (t_1^+,t_1,t_1^{++},t')
  \\[1mm]
  &
  -2
  [G_D]_{uq}(t_1,t')
  [C_2^D]_{tp,sr}
  (t_1,t,t_1^+,t_1^{++})
  \\[1mm]
  &
  -4
  [G_D]_{tr}(t_1,t_1^+)
  [C_2^D]_{up,qs}
  (t_1,t,t',t_1^+)
  \\[1mm]
  &
  -4
  [G_D]_{pr}(t,t_1)
  [G_D]_{us}(t_1,t_1^+)
  [G_D]_{tq}(t_1,t')
  \Big\}
  \end{aligned}
  \label{eq:G1-twobody}
\end{equation}

Here
\begin{equation}
  t_1^{+++}
  >
  t_1^{++}
  >
  t_1^+
  >
  t_1
\end{equation}
denotes the conventional infinitesimal time ordering at the
interaction vertex.

The numerical coefficients in
Eq.~\eqref{eq:G1-twobody} are important.  The two terms containing one
one-particle propagator and one connected two-particle propagator carry
the factor
\begin{equation}
  \frac{i}{4}(-2)
  =
  -\frac{i}{2}
\end{equation}
which will generate $\Sigma_1^{12}$ and $\Sigma_1^{21}$.

Likewise,
\begin{equation}
  \frac{i}{4}(-4)
  =
  -i
\end{equation}
which generates the Hartree--Fock-like one-body contribution and the
density-dependent contribution to $\Sigma_1^{22}$.

Finally, the fully connected three-body term retains the prefactor
\begin{equation}
  \frac{i}{4}
\end{equation}


\subsubsection{Regrouping the first-order Green's function}

Define the zeroth-order one-body density matrix by
\begin{equation}
  \gamma_{pq}
  =
  \braket{
    \hat p^\dagger\hat q
  }_D
  =
  (-i)
  [G_D]_{qp}(t,t^+)
  \label{eq:reference-density}
\end{equation}

The effective static one-body perturbation is
\begin{equation}
  u^{\mathrm{eff}}_{rs}
  =
  u_{rs}
  +
  \bar v_{rs,tu}\gamma_{tu}
  \label{eq:u-effective}
\end{equation}

Combining Eqs.~\eqref{eq:G1-onebody} and
\eqref{eq:G1-twobody}, and using the antisymmetry of $\bar v$, the
first-order Green's function can be regrouped as
\begin{equation}
  \begin{aligned}
  [G_1]_{pq}(t,t')
  =
  \int dt_1
  \Big\{
  &
  [G_D]_{pr}(t,t_1)
  u^{\mathrm{eff}}_{rs}
  [G_D]_{sq}(t_1,t')
  \\[1mm]
  &
  +
  [G_D]_{pr}(t,t_1)
  \left[
  -\frac{i}{2}
  \bar v_{rt,su}
  [C_2^D]_{ut,sq}
  (t_1^+,t_1,t_1^{++},t')
  \right]
  \\[1mm]
  &
  +
  \left[
  -\frac{i}{2}
  [C_2^D]_{pt,rs}
  (t,t_1,t_1^{++},t_1^+)
  \bar v_{rt,su}
  \right]
  [G_D]_{uq}(t_1,t')
  \\[1mm]
  &
  -
  u^{\mathrm{eff}}_{rs}
  [C_2^D]_{ps,qr}
  (t,t_1,t',t_1^+)
  \\[1mm]
  &
  +
  \frac{i}{4}
  \bar v_{rt,su}
  [C_3^D]_{ptu,qrs}
  (t,t_1,t_1^+,t',t_1^{+++},t_1^{++})
  \Big\}
  \end{aligned}
  \label{eq:G1-regrouped}
\end{equation}

Equation~\eqref{eq:G1-regrouped} has exactly the four structures
appearing in the first-order generalized Dyson equation,
Eq.~\eqref{eq:first-order-generalized-dyson}.


\subsubsection{Identification of the four self-energy blocks}

Comparison of Eqs.~\eqref{eq:first-order-generalized-dyson} and
\eqref{eq:G1-regrouped} gives the four first-order generalized
self-energy blocks.

\paragraph{Physical block.}

The $11$ block is local in time,
\begin{equation}
  [\Sigma_1^{11}]_{rs}(t,t')
  =
  \delta(t-t')
  u^{\mathrm{eff}}_{rs}
  \label{eq:sigma11-first}
\end{equation}
with
\begin{equation}
  u^{\mathrm{eff}}_{rs}
  =
  u_{rs}
  +
  \bar v_{rs,tu}\gamma_{tu}
\end{equation}


\paragraph{Right mixed block.}

The $12$ block is
\begin{equation}
  [\Sigma_1^{12}]_{rq}(t,t')
  =
  -\frac{i}{2}
  \bar v_{rt,su}
  [C_2^D]_{ut,sq}
  (t^+,t,t^{++},t')
  \label{eq:sigma12-first}
\end{equation}


\paragraph{Left mixed block.}

The $21$ block is
\begin{equation}
  [\Sigma_1^{21}]_{pu}(t,t')
  =
  -\frac{i}{2}
  [C_2^D]_{pt,rs}
  (t,t',t'^{++},t'^+)
  \bar v_{rt,su}
  \label{eq:sigma21-first}
\end{equation}


\paragraph{Auxiliary block.}

Finally,
\begin{equation}
  \begin{aligned}
  [\Sigma_1^{22}]_{pq}(t,t')
  =
  \int dt_1
  \Big[
  &
  -
  u^{\mathrm{eff}}_{rs}
  [C_2^D]_{ps,qr}
  (t,t_1,t',t_1^+)
  \\
  &
  +
  \frac{i}{4}
  \bar v_{rt,su}
  [C_3^D]_{ptu,qrs}
  (t,t_1,t_1^+,t',t_1^{+++},t_1^{++})
  \Big]
  \end{aligned}
  \label{eq:sigma22-first}
\end{equation}

Equations~\eqref{eq:sigma11-first}--\eqref{eq:sigma22-first}
coincide with Eqs.~(20)--(23) of Ref.~\cite{WangFangLi2026}.


\subsubsection{Origin of the numerical prefactors}

It is useful to emphasize the origin of the characteristic prefactors.

After substitution of the three-body cumulant decomposition into the
antisymmetrized two-electron perturbation and combination of
permutation-equivalent contractions, the terms generating the mixed blocks
acquire a factor $-2$, whereas the density-dependent contractions acquire a
factor $-4$.  Combined with the overall prefactor $i/4$, these give
\begin{equation}
  \frac{i}{4}\times(-2)
  =
  -\frac{i}{2}
  \qquad
  \frac{i}{4}\times(-4)
  =
  -i
\end{equation}
The genuinely connected three-body contribution retains the prefactor
\begin{equation}
  +\frac{i}{4}
\end{equation}

These factors provide a useful independent check of the enlarged-source
conventions introduced in Sec.~\ref{sec:dyall-hedin}.


\subsubsection{Specialization to the Dyall Hamiltonian}

We now exploit the special structure of the Dyall reference.

The Dyall zeroth-order eigenstates factorize into inactive and active
parts.  The inactive sector is Gaussian, while the active sector
contains the non-Gaussian correlations.  Consequently, any connected
zeroth-order Green's function of order two or higher vanishes if one
of its orbital indices is inactive:
\begin{equation}
  C_n^D
  \neq 0
  \quad\Longrightarrow\quad
  \text{all orbital indices of }C_n^D
  \text{ are active}
  \qquad
  n\ge2
  \label{eq:cumulant-active-support}
\end{equation}

Furthermore, the residual two-electron interaction vanishes when all
four indices are active,
\begin{equation}
  v^R_{xy,zw}=0
  \label{eq:residual-all-active-zero}
\end{equation}


\subsubsection{Dyall simplification of the physical block}

For the Dyall partition,
\begin{equation}
  [\Sigma_1^{11}]_{rs}
  =
  u_{rs}
  +
  \bar v_{rs,pq}\gamma_{pq}
\end{equation}
is frequency independent and Hermitian.

Partitioning the inactive orbitals into core indices $i,j,\ldots$ and
virtual indices $a,b,\ldots$, the independent matrix elements are
\begin{subequations}
  \label{eq:dyall-sigma11-blocks}
  \begin{align}
    [\Sigma_1^{11}]_{ij}
    &=
    0
    \\
    [\Sigma_1^{11}]_{aj}
    &=
    h_{aj}
    +
    \bar v_{aj,pq}\gamma_{pq}
    \\
    [\Sigma_1^{11}]_{xj}
    &=
    h_{xj}
    +
    \bar v_{xj,pq}\gamma_{pq}
    \\
    [\Sigma_1^{11}]_{ab}
    &=
    0
    \\
    [\Sigma_1^{11}]_{xb}
    &=
    h_{xb}
    +
    \bar v_{xb,pq}\gamma_{pq}
    \\
    [\Sigma_1^{11}]_{xy}
    &=
    0
  \end{align}
\end{subequations}

The remaining blocks follow by Hermiticity.

In particular,
\begin{equation}
  [\Sigma_1^{11}]_{xy}=0
  \label{eq:active-active-sigma11-zero}
\end{equation}
because the active--active mean-field contribution already contained
in the Dyall Hamiltonian exactly cancels the corresponding part of the
residual one-body perturbation.


\subsubsection{Vanishing of \texorpdfstring{$\Sigma_1^{22}$}{Sigma22} for the Dyall reference}

Equation~\eqref{eq:sigma22-first} contains two terms.

For the first term to be nonzero,
\begin{equation}
  [C_2^D]_{ps,qr}
  \neq0
\end{equation}
all four indices $p,s,q,r$ must be active by
Eq.~\eqref{eq:cumulant-active-support}.  Hence
\begin{equation}
  u^{\mathrm{eff}}_{rs}
  =
  [\Sigma_1^{11}]_{rs}
  =
  0
\end{equation}
by Eq.~\eqref{eq:active-active-sigma11-zero}.

For the second term,
\begin{equation}
  [C_3^D]_{ptu,qrs}
  \neq0
\end{equation}
requires all six orbital indices to be active.  In particular,
$r,t,s,u$ are active, and therefore
\begin{equation}
  \bar v_{rt,su}=0
\end{equation}
by Eq.~\eqref{eq:residual-all-active-zero}.

In the diagrammatic formulation of Ref.~\cite{WangFangLi2026}, the two
contributions involving the one-body and mean-field pieces cancel upon use
of the Dyall relation that leads to
$u^{\mathrm{eff}}_{xy}=[\Sigma_1^{11}]_{xy}=0$.  In the regrouped expression
of Eq.~\eqref{eq:sigma22-first}, this cancellation is already encoded in the
vanishing of its first term.  The remaining connected three-body term
vanishes separately because the residual interaction has no all-active
matrix elements, $v^R_{xy,zw}=0$.  Therefore,
\begin{equation}
  \Sigma_1^{22}=0
  \label{eq:dyall-sigma22-zero}
\end{equation}


\subsubsection{Triangular structure of the mixed blocks}

Consider first
\begin{equation}
  [\Sigma_1^{12}]_{rq}
  =
  -\frac{i}{2}
  \bar v_{rt,su}
  [C_2^D]_{ut,sq}
\end{equation}

A nonzero connected two-body Green's function requires
\begin{equation}
  u,t,s,q\in A
\end{equation}
Thus $q$ must be active.  If $r$ were also active, all four indices of
$\bar v_{rt,su}$ would be active, so that the residual interaction
would vanish.  Hence $r$ must be inactive.

Therefore
\begin{equation}
  \Sigma_1^{12}
  =
  \begin{pmatrix}
    0 & \Sigma_1^{12,IA}
    \\
    0 & 0
  \end{pmatrix}
  \label{eq:sigma12-triangular}
\end{equation}

Similarly,
\begin{equation}
  [\Sigma_1^{21}]_{pu}
  =
  -\frac{i}{2}
  [C_2^D]_{pt,rs}
  \bar v_{rt,su}
\end{equation}
requires
\begin{equation}
  p,t,r,s\in A
\end{equation}
Thus $p$ is active.  If $u$ were also active, the residual interaction
would again contain four active indices and vanish.  Hence $u$ must be
inactive.

Consequently,
\begin{equation}
  \Sigma_1^{21}
  =
  \begin{pmatrix}
    0 & 0
    \\
    \Sigma_1^{21,AI} & 0
  \end{pmatrix}
  \label{eq:sigma21-triangular}
\end{equation}

Equations~\eqref{eq:sigma12-triangular} and
\eqref{eq:sigma21-triangular} reproduce the block structure given in
Eq.~(S14) of Ref.~\cite{WangFangLi2026}.


\subsubsection{Spectral representation of the mixed blocks}

The triangular structure permits particularly simple spectral
representations.  The expressions in this subsection are the time-ordered
(Feynman) forms of Ref.~\cite{WangFangLi2026}; accordingly, addition and
removal poles carry opposite infinitesimal prescriptions.  The conversion to
a retarded representation, in which all signed poles carry the same
$+i0^+$ prescription, is performed explicitly later in
Sec.~\ref{sec:hermitian-completion}.

Let
\begin{equation}
  \hat H_D^{A}
  \ket{\Xi_{\mu,N_a}^{A}}
  =
  E_{\mu,N_a}^{A}
  \ket{\Xi_{\mu,N_a}^{A}}
\end{equation}
and define the active charged excitation energies
\begin{equation}
  \omega_{\mu,N_a\pm1}^{A}
  =
  E_{\mu,N_a\pm1}^{A}
  -
  E_{0,N_a}^{A}
\end{equation}

Insertion of complete sets of $N_a+1$ and $N_a-1$ active-space
eigenstates into Eq.~\eqref{eq:sigma12-first} gives
\begin{equation}
  \begin{aligned}
  [\Sigma_1^{12,IA}]_{Px}(\omega)
  ={}&
  \frac{1}{2}
  \bar v_{Pz,yw}
  \sum_{\mu}
  \frac{
    \braket{
      \Xi_{0,N_a}^{A}
      |
      \hat y^\dagger\hat w\hat z
      |
      \Xi_{\mu,N_a+1}^{A}
    }
    \braket{
      \Xi_{\mu,N_a+1}^{A}
      |
      \hat x^\dagger
      |
      \Xi_{0,N_a}^{A}
    }
  }{
    \omega
    -
    \omega_{\mu,N_a+1}^{A}
    +
    i0^+
  }
  \\[1mm]
  &+
  \frac{1}{2}
  \bar v_{Pz,yw}
  \sum_{\mu}
  \frac{
    \braket{
      \Xi_{0,N_a}^{A}
      |
      \hat x^\dagger
      |
      \Xi_{\mu,N_a-1}^{A}
    }
    \braket{
      \Xi_{\mu,N_a-1}^{A}
      |
      \hat y^\dagger\hat w\hat z
      |
      \Xi_{0,N_a}^{A}
    }
  }{
    \omega
    +
    \omega_{\mu,N_a-1}^{A}
    -
    i0^+
  }
  \\[1mm]
  &-
  \bar v_{Pz,yw}
  \braket{
    \hat y^\dagger\hat w
  }_D
  [G_D^A]_{zx}(\omega)
  \end{aligned}
  \label{eq:sigma12-spectral}
\end{equation}

Likewise,
\begin{equation}
  \begin{aligned}
  [\Sigma_1^{21,AI}]_{xP}(\omega)
  ={}&
  -\frac{1}{2}
  \bar v_{zP,wy}
  \sum_{\mu}
  \frac{
    \braket{
      \Xi_{0,N_a}^{A}
      |
      \hat x
      |
      \Xi_{\mu,N_a+1}^{A}
    }
    \braket{
      \Xi_{\mu,N_a+1}^{A}
      |
      \hat z^\dagger\hat w^\dagger\hat y
      |
      \Xi_{0,N_a}^{A}
    }
  }{
    \omega
    -
    \omega_{\mu,N_a+1}^{A}
    +
    i0^+
  }
  \\[1mm]
  &-
  \frac{1}{2}
  \bar v_{zP,wy}
  \sum_{\mu}
  \frac{
    \braket{
      \Xi_{\mu,N_a-1}^{A}
      |
      \hat x
      |
      \Xi_{0,N_a}^{A}
    }
    \braket{
      \Xi_{0,N_a}^{A}
      |
      \hat z^\dagger\hat w^\dagger\hat y
      |
      \Xi_{\mu,N_a-1}^{A}
    }
  }{
    \omega
    +
    \omega_{\mu,N_a-1}^{A}
    -
    i0^+
  }
  \\[1mm]
  &+
  \bar v_{zP,wy}
  \braket{
    \hat w^\dagger\hat y
  }_D
  [G_D^A]_{xz}(\omega)
  \end{aligned}
  \label{eq:sigma21-spectral}
\end{equation}

These are Eqs.~(S15) and (S16) of
Ref.~\cite{WangFangLi2026}.


\subsubsection{Interpretation of the connected subtraction}

The last terms of Eqs.~\eqref{eq:sigma12-spectral} and
\eqref{eq:sigma21-spectral} arise from the subtraction of the
disconnected part of the two-body Green's function.

For example, the attachment part of
Eq.~\eqref{eq:sigma12-spectral} contains the combination
\begin{equation}
  \frac{1}{2}
  \braket{
    \Xi_0
    |
    \hat y^\dagger\hat w\hat z
    |
    \Xi_\mu^{+}
  }
  -
  \braket{
    \hat y^\dagger\hat w
  }
  \braket{
    \Xi_0
    |
    \hat z
    |
    \Xi_\mu^{+}
  }
\end{equation}

Thus the mixed self-energy is driven by the genuinely connected
correlation between the one-particle propagation and the active-space
density fluctuation.  It vanishes identically for a Gaussian
zeroth-order reference.
%%%%%%%%%%%%%%%%%%%%%%%%%
\subsection{The existing MR-\texorpdfstring{$GW$}{GW} approximation as a Dyall-Hedin truncation}
\label{sec:mr-gw-from-hedin}
%%%%%%%%%%%%%%%%%%%%%%%%%

\Established The practical MR-$GW$ approximation in this subsection is
the 2026 Wang--Fang--Li scheme.  The new element is its interpretation as a
channel-selective truncation of the Dyall-Hedin hierarchy derived above.


Having established the exact Dyall-Hedin hierarchy and verified its
first-order expansion, we can now identify precisely the approximations
that lead to the practical MR-$GW$ method of Wang, Fang, and Li
\cite{WangFangLi2026}.

This connection is important for two reasons.  First, it provides a
functional interpretation of an approximation originally introduced
through generalized diagrammatics and Hall's generalized Dyson
equation.  Second, it identifies unambiguously which pieces of the exact
Dyall-Hedin hierarchy have been retained and which have been discarded,
thereby defining a systematic path beyond the existing method.


\subsubsection{One-shot Dyall-reference approximation}

The exact theory contains the dressed enlarged Green's function
$\mathbb G$, the enlarged vertex $\overline{\boldsymbol\Gamma}$, the
polarization $P_R$, and the residual screened interaction $W_R$.

The practical MR-$GW$ approximation is one-shot.  All quantities used to
construct the polarization and self-energy are therefore evaluated from
the fixed Dyall reference,
\begin{equation}
  G \longrightarrow G_D
\end{equation}
inside the self-energy construction, while the final interacting Green's
function is obtained from a Dyson equation.

This is the direct multireference analogue of the $G_0W_0$ strategy.


\subsubsection{Reference polarization}

The exact Dyall-Hedin polarization is
\begin{equation}
  P_R
  =
  -i
  \left[
    \mathbb G
    \overline{\boldsymbol\Gamma}
    \mathbb G
  \right]_{11}^{ph}
  \label{eq:exact-P-task3}
\end{equation}

At the fixed-reference level, no residual-interaction dressing of the
response is retained before constructing the RPA screening.  Thus
\begin{equation}
  P_R
  \longrightarrow
  P_D
  =
  -i
  \left[
    \mathbb G_D
    \overline{\boldsymbol\Gamma}_D
    \mathbb G_D
  \right]_{11}^{ph}
  \label{eq:PR-to-PD}
\end{equation}

Using the exact Dyall Schwinger relation,
\begin{equation}
  \left[
    \mathbb G_D
    \overline{\boldsymbol\Gamma}_D
    \mathbb G_D
  \right]_{11}
  =
  G_DG_D-C_2^{D,ph}
\end{equation}
one obtains
\begin{equation}
  P_D
  =
  -i
  \left(
    G_DG_D-C_2^{D,ph}
  \right)
  \label{eq:PD-task3}
\end{equation}

Therefore the reference polarization is \emph{not} an independent-particle
bubble.

Equivalently, introducing the effective physical Dyall vertex,
\begin{equation}
  G_D\Gamma_DG_D
  =
  G_DG_D-C_2^{D,ph}
\end{equation}
one may write
\begin{equation}
  P_D
  =
  -iG_D\Gamma_DG_D
  \label{eq:PD-GammaD}
\end{equation}

The nontrivial reference vertex $\Gamma_D$ is thus retained exactly in
the screening channel.


\subsubsection{Equivalence with the MR-RPA irreducible polarizability}

The Lehmann representation of Eq.~\eqref{eq:PD-task3} is
\begin{equation}
  \begin{aligned}
  [P_D]_{pr,qs}(\omega)
  ={}&
  \sum_{\mu>0}
  \frac{
    \braket{
      \Phi_0
      |
      \hat p^\dagger\hat r
      |
      \Phi_\mu
    }
    \braket{
      \Phi_\mu
      |
      \hat q^\dagger\hat s
      |
      \Phi_0
    }
  }{
    \omega-\omega_\mu+i0^+
  }
  \\
  &-
  \sum_{\mu>0}
  \frac{
    \braket{
      \Phi_0
      |
      \hat q^\dagger\hat s
      |
      \Phi_\mu
    }
    \braket{
      \Phi_\mu
      |
      \hat p^\dagger\hat r
      |
      \Phi_0
    }
  }{
    \omega+\omega_\mu-i0^+
  }
  \end{aligned}
  \label{eq:PD-lehmann}
\end{equation}

Equation~\eqref{eq:PD-lehmann} is precisely the quantity denoted
$\Pi_0(\omega)$ in the MR-RPA construction of
Ref.~\cite{WangFangLi2026}.  Hence,
\begin{equation}
  \Pi_0^{\mathrm{MR-RPA}}
  =
  P_D
  \label{eq:Pi0-equals-PD}
\end{equation}

This equality provides the first direct connection between the
Dyall-Hedin hierarchy and the published MR-$GW$ construction.


\subsubsection{Residual screened interaction}

The exact screened residual interaction satisfies
\begin{equation}
  W_R
  =
  v_R
  +
  v_RP_RW_R
\end{equation}

Using the fixed Dyall polarization,
Eq.~\eqref{eq:PR-to-PD}, gives
\begin{equation}
  W_D
  =
  v_R
  +
  v_RP_DW_D
  \label{eq:WD-task3}
\end{equation}

This is exactly the screened interaction employed in the practical
MR-$GW$ method: the full Coulomb interaction of conventional $GW$ is
replaced by the residual two-electron interaction, while the
independent-particle RPA polarizability is replaced by the correlated
Dyall response.

The active--active Coulomb interaction must not be screened again,
because it is already contained nonperturbatively in $H_D$.


\subsubsection{An important distinction from conventional \texorpdfstring{$GW$}{GW}}

In conventional Hedin theory,
\begin{equation}
  P=-iGG\Gamma
\end{equation}
and
\begin{equation}
  \Sigma=iGW\Gamma
\end{equation}
so that setting
\begin{equation}
  \Gamma\rightarrow1
\end{equation}
simultaneously generates the RPA polarization and the $GW$ self-energy.

The multireference construction is more subtle.

If one were to set
\begin{equation}
  \Gamma_D\rightarrow1
\end{equation}
in the Dyall reference response, one would obtain
\begin{equation}
  P_D
  \longrightarrow
  -iG_DG_D
\end{equation}
and the connected contribution
\begin{equation}
  iC_2^{D,ph}
\end{equation}
would be lost.

This is \emph{not} the MR-RPA polarizability used by
Ref.~\cite{WangFangLi2026}.

Conversely, inserting the full $\Gamma_D$ explicitly into the
self-energy would generate an approximation of the form
\begin{equation}
  iG_DW_D\Gamma_D
\end{equation}
which is also \emph{not} the published MR-$GW$ approximation.

The existing method therefore cannot be characterized by the single
replacement
\begin{equation}
  \Gamma
  \rightarrow
  \Gamma_D
\end{equation}

Nor can it be characterized solely by
\begin{equation}
  \Delta\Gamma_R=0
\end{equation}

Instead, it constitutes a \emph{channel-selective} truncation of the
Dyall-Hedin hierarchy.


\subsubsection{Self-energy truncation}

The practical approximation retains, in the physical $11$ block, the
same self-energy topology as conventional one-shot $GW$.

From Sec.~\ref{sec:first-order-blocks}, the static first-order physical block is
\begin{equation}
  [\Sigma_1^{11}]_{pq}
  =
  u_{pq}
  +
  \bar v_{pq,rs}
  \gamma_{rs}^D
  \label{eq:sigma11-static-task3}
\end{equation}

The correlation contribution is obtained from the usual $GW$
convolution, but using the interacting Dyall Green's function and the
MR-RPA screened residual interaction,
\begin{equation}
  \Sigma_{pq}^{c,\mathrm{MR-}GW}(\omega)
  =
  \frac{i}{2\pi}
  \int d\omega'
  \
  [G_D]_{rs}(\omega+\omega')
  [w_D^c]_{pr,sq}(\omega')
  \label{eq:MRGW-convolution}
\end{equation}

The resulting physical self-energy is therefore
\begin{equation}
  \Sigma^{\mathrm{MR-}GW}
  =
  \Sigma_1^{11}
  +
  \Sigma^{c,\mathrm{MR-}GW}
  \label{eq:MRGW-self-energy}
\end{equation}

Equation~\eqref{eq:MRGW-self-energy} is exactly the quantity defined in
Eq.~(S41) of Ref.~\cite{WangFangLi2026}.


\subsubsection{Spectral form}

Using the Lehmann representation of $G_D$ and the MR-RPA spectral
representation
\begin{equation}
  [w_D^c]_{pr,qs}(\omega)
  =
  \sum_{I>0}
  \left[
    \frac{
      M_{pr,I}M_{qs,I}^{*}
    }{
      \omega-\Omega_I+i0^+
    }
    -
    \frac{
      M_{pr,I}^{*}M_{qs,I}
    }{
      \omega+\Omega_I-i0^+
    }
  \right]
\end{equation}
one obtains
\begin{equation}
  \begin{aligned}
  \Sigma_{pq}^{\mathrm{MR-}GW}(\omega)
  ={}&
  [\Sigma_1^{11}]_{pq}
  \\
  &+
  \sum_{I>0}
  \sum_{\mu}
  \frac{
    \braket{
      \Phi_0
      |
      \hat r
      |
      \Phi_\mu^{N+1}
    }
    \braket{
      \Phi_\mu^{N+1}
      |
      \hat s^\dagger
      |
      \Phi_0
    }
    M_{pr,I}^{*}
    M_{sq,I}
  }{
    \omega
    -
    \Omega_I
    -
    \omega_\mu^{N+1}
    +
    i0^+
  }
  \\
  &+
  \sum_{I>0}
  \sum_{\mu}
  \frac{
    \braket{
      \Phi_0
      |
      \hat s^\dagger
      |
      \Phi_\mu^{N-1}
    }
    \braket{
      \Phi_\mu^{N-1}
      |
      \hat r
      |
      \Phi_0
    }
    M_{pr,I}
    M_{sq,I}^{*}
  }{
    \omega
    +
    \Omega_I
    +
    \omega_\mu^{N-1}
    -
    i0^+
  }
  \end{aligned}
  \label{eq:MRGW-spectral-task3}
\end{equation}

This is the general spectral expression given in
Ref.~\cite{WangFangLi2026}.


\subsubsection{Truncation of the remaining Hall blocks}

The exact generalized Dyson equation contains
\begin{equation}
  \Sigma^{11}
  \qquad
  \Sigma^{12}
  \qquad
  \Sigma^{21}
  \qquad
  \Sigma^{22}
\end{equation}

The practical MR-$GW$ approximation sets
\begin{equation}
  \Sigma^{12}
  =
  \Sigma^{21}
  =
  0
  \label{eq:MRGW-mixed-zero}
\end{equation}

For the Dyall Hamiltonian we have already established that
\begin{equation}
  \Sigma_1^{22}=0
\end{equation}
The practical approximation additionally neglects all higher-order
contributions to this block,
\begin{equation}
  \Sigma^{22}=0
  \label{eq:MRGW-22-zero}
\end{equation}

Consequently,
\begin{equation}
  M
  =
  G_D
\end{equation}

Hall's generalized Dyson equation therefore collapses to
\begin{equation}
  G
  =
  G_D
  +
  G_D
  \Sigma^{\mathrm{MR-}GW}
  G
  \label{eq:MRGW-Dyson-task3}
\end{equation}

Equivalently,
\begin{equation}
  G(\omega)
  =
  \left[
    G_D^{-1}(\omega)
    -
    \Sigma^{\mathrm{MR-}GW}(\omega)
  \right]^{-1}
  \label{eq:MRGW-inverse-task3}
\end{equation}

This is precisely the working Dyson equation of
Ref.~\cite{WangFangLi2026}.


\subsubsection{Interpretation as a channel-selective vertex truncation}

The complete mapping can now be stated precisely.

In the \emph{screening channel}, the full Dyall reference response is
retained:
\begin{equation}
  G_D\Gamma_DG_D
  =
  G_DG_D-C_2^{D,ph}
\end{equation}
Thus
\begin{equation}
  \Gamma_{\mathrm{screen}}
  =
  \Gamma_D
  \label{eq:Gamma-screen}
\end{equation}

In the \emph{physical self-energy channel}, only the ordinary $GW$
self-energy topology is retained:
\begin{equation}
  \Sigma^{11}
  \longrightarrow
  \Sigma_1^{11}
  +
  iG_DW_D^c
  \label{eq:Gamma-sigma11}
\end{equation}
where the product in the second term denotes the usual $GW$
frequency convolution.

Thus no explicit reference vertex $\Gamma_D$ is attached to the
self-energy $GW$ diagram.

Finally, in the \emph{mixed and auxiliary channels},
\begin{equation}
  \Sigma^{12}
  =
  \Sigma^{21}
  =
  \Sigma^{22}
  =
  0
  \label{eq:mixed-aux-truncation}
\end{equation}

The practical method is therefore most accurately described as
\begin{equation}
  \begin{array}{rcl}
  \text{screening channel:}
  &&
  \Gamma_D
  \text{ retained}
  \\[1mm]
  \text{$11$ self-energy channel:}
  &&
  \text{bare $GW$ vertex retained}
  \\[1mm]
  \text{mixed/auxiliary channels:}
  &&
  \text{discarded}
  \end{array}
  \label{eq:channel-selective-summary}
\end{equation}

This selective treatment is the precise Dyall-Hedin interpretation of
the existing MR-$GW$ approximation.


\subsubsection{Relation to the statement
\texorpdfstring{$\Delta\Gamma_R=0$}{Delta Gamma R = 0}}

It is still useful to write formally
\begin{equation}
  \Gamma
  =
  \Gamma_D
  +
  \Delta\Gamma_R
\end{equation}
where $\Delta\Gamma_R$ denotes vertex corrections induced by the
residual interaction.

The approximation
\begin{equation}
  \Delta\Gamma_R=0
\end{equation}
correctly expresses the fact that no additional residual-interaction
vertex equation is solved when constructing $P_D$ and $W_D$.

However, it is \emph{not sufficient} to define the practical MR-$GW$
self-energy, because retaining $\Gamma_D$ explicitly in the
self-energy would produce a $GW\Gamma_D$-type theory.

The full definition of the published MR-$GW$ approximation therefore
requires both
\begin{equation}
  \Delta\Gamma_R=0
  \quad\text{in the reference screening problem}
\end{equation}
and the additional projected self-energy truncations
\begin{equation}
  \Sigma^{11}
  \rightarrow
  \Sigma_1^{11}
  +
  iG_DW_D^c
\end{equation}
together with
\begin{equation}
  \Sigma^{12}
  =
  \Sigma^{21}
  =
  \Sigma^{22}
  =
  0
\end{equation}


\subsubsection{Why the construction is internally consistent}

Although the vertex is treated differently in the screening and
self-energy channels, this does not imply double counting.

The interaction contained completely within the active space belongs to
$H_D$ and is therefore absent from the residual interaction $v_R$.
The strong active-space correlation consequently enters through
\begin{equation}
  G_D
  \quad\text{and}\quad
  P_D
\end{equation}
whereas $W_D$ screens only the residual interaction.

The final self-energy therefore separates naturally into
\begin{equation}
  \text{nonperturbative active-space correlation}
  +
  \text{screened residual correlation}
\end{equation}

This is also evident if one introduces a noninteracting propagator
$g_0$ and defines
\begin{equation}
  \Sigma_D
  =
  g_0^{-1}
  -
  G_D^{-1}
\end{equation}
Then
\begin{equation}
  G^{-1}
  =
  g_0^{-1}
  -
  \left(
    \Sigma_D
    +
    \Sigma^{\mathrm{MR-}GW}
  \right)
\end{equation}
which is the standard-Dyson representation discussed in
Ref.~\cite{WangFangLi2026}.


\subsubsection{MR-\texorpdfstring{$GW$}{GW} is not first-order exact for a
correlated reference}

One important consequence of the above identification should be stated
explicitly.

From Sec.~\ref{sec:first-order-blocks}, the exact first-order Green's function is
\begin{equation}
  G_1
  =
  G_D\Sigma_1^{11}G_D
  +
  G_D\Sigma_1^{12}
  +
  \Sigma_1^{21}G_D
  +
  \Sigma_1^{22}
\end{equation}

For the Dyall Hamiltonian,
\begin{equation}
  \Sigma_1^{22}=0
\end{equation}
but in general
\begin{equation}
  \Sigma_1^{12}\neq0
  \qquad
  \Sigma_1^{21}\neq0
\end{equation}

The practical MR-$GW$ approximation discards these terms.

Therefore,
\begin{equation}
  \text{MR-}GW
  \text{ does not reproduce the complete first-order expansion}
\end{equation}
for a genuinely correlated Dyall reference.

This is not an inconsistency: MR-$GW$ is a deliberately selected
screened approximation rather than a finite-order perturbation theory.

However, this observation identifies immediately the lowest-cost
information discarded by the existing method:
\begin{equation}
  \Sigma^{12}
  \quad\text{and}\quad
  \Sigma^{21}
\end{equation}

Restoring these terms in a screened and analytically consistent manner
therefore provides the natural starting point for a systematic
extension of MR-$GW$.


\subsubsection{Objective I in compact form}

The derivation establishes the following sequence:
\begin{equation}
  \begin{aligned}
  &\text{exact Dyall-Hedin hierarchy}
  \\
  &\quad
  \xrightarrow{
    \substack{
    \text{freeze the response}\\
    \text{at the Dyall reference}
    }
  }
  P_D
  =
  -iG_D\Gamma_DG_D
  \\
  &\quad
  \xrightarrow{
    \text{MR-RPA}
  }
  W_D
  =
  v_R+v_RP_DW_D
  \\
  &\quad
  \xrightarrow{
    \substack{
    \text{retain only the ordinary}\\
    GW\text{ topology in }\Sigma^{11}
    }
  }
  \Sigma^{\mathrm{MR-}GW}
  =
  \Sigma_1^{11}
  +
  iG_DW_D^c
  \\
  &\quad
  \xrightarrow{
    \Sigma^{12}
    =
    \Sigma^{21}
    =
    \Sigma^{22}
    =
    0
  }
  G
  =
  G_D
  +
  G_D
  \Sigma^{\mathrm{MR-}GW}
  G
  \end{aligned}
  \label{eq:objective-I-summary}
\end{equation}

Thus the MR-$GW$ method of Ref.~\cite{WangFangLi2026} can be
identified as a specific \emph{channel-selective vertex and block
truncation} of the Dyall-Hedin hierarchy.

This completes the first objective of the present development and,
at the same time, identifies the mixed blocks
$\Sigma^{12}$ and $\Sigma^{21}$ as the most natural targets for the
first practical extension beyond the existing MR-$GW$ approximation.
%%%%%%%%%%%%%%%%%%%%%%%%%
\section{Objective II: Screened mixed blocks and a Hermitian PSD completion}
\label{sec:objective-II}
%%%%%%%%%%%%%%%%%%%%%%%%%

The exact and first-order analyses of Objective~I identify
$\Sigma^{12}$ and $\Sigma^{21}$ as the lowest-cost information discarded by
the practical MR-$GW$ approximation.  Objective~II restores these blocks in a
controlled sequence: first by defining a parent-consistent static screened
kernel, then by embedding the screened residues in a Hermitian pole-space
representation, and finally by analyzing spectral moments and computational
cost.
%%%%%%%%%%%%%%%%%%%%%%%%%
\subsection{A parent-consistent static kernel for the mixed Hall sector}
\label{sec:static-mixed-kernel}
%%%%%%%%%%%%%%%%%%%%%%%%%

\Proposed The static kernel introduced below is our first approximation
beyond the published MR-$GW$ scheme.  It is not an exact identity: it follows
from the parent MR-$GW$ self-energy after the frozen-$W$ and static-screening
approximations.


The first extension beyond the MR-$GW$ approximation should restore
information contained in the mixed Hall blocks
\begin{equation}
  \Sigma^{12}
  \qquad\text{and}\qquad
  \Sigma^{21}
\end{equation}
while preserving the favorable analytical properties of the existing
MR-$GW$ approximation.

At first order, these blocks are generated by the antisymmetrized
residual interaction,
\begin{equation}
  \bar v^R_{pr,qs}
  =
  v^R_{pr,qs}
  -
  v^R_{ps,qr}
  \label{eq:bare-antisym-residual}
\end{equation}
through
\begin{equation}
  [\Sigma^{12}_{1}]_{rq}
  =
  -\frac{i}{2}
  \bar v^R_{rt,su}
  [C_2^D]_{ut,sq}
\end{equation}
and
\begin{equation}
  [\Sigma^{21}_{1}]_{pu}
  =
  -\frac{i}{2}
  [C_2^D]_{pt,rs}
  \bar v^R_{rt,su}
\end{equation}

The question addressed in this section is how the bare interaction in
these expressions should be screened in the cheapest approximation
consistent with the parent MR-$GW$ theory.


\subsubsection{Requirements on the screened mixed kernel}

We impose the following requirements.

\paragraph{Consistency with the parent MR-\texorpdfstring{$GW$}{GW}
approximation.}

The kernel should follow from the same physical self-energy used to
construct MR-$GW$, rather than being introduced independently.

\paragraph{Correct bare limit.}

When screening is removed,
\begin{equation}
  W_D
  \longrightarrow
  v_R
\end{equation}
the kernel must reduce to the antisymmetrized residual interaction,
\begin{equation}
  \mathcal I_D^{\rm mix}
  \longrightarrow
  \bar v_R
\end{equation}

\paragraph{No additional dynamical vertex problem.}

The first extension should not require the solution of a new
two-frequency Bethe--Salpeter equation.

\paragraph{Compatibility with a Hermitian/PSD completion.}

The static kernel should possess the Hermiticity properties required to
construct conjugate left and right mixed couplings.

\paragraph{Correct multireference limits.}

The mixed contribution must vanish in the Gaussian-reference limit and
in the full-active-space limit.


\subsubsection{The parent MR-\texorpdfstring{$GW$}{GW} self-energy}

It is useful to reorganize the physical $11$ self-energy of the
existing MR-$GW$ approximation into a Hartree contribution and a
$GW$ exchange-correlation contribution.

Suppressing the one-body residual term $u$, which is independent of the
Green's function, we may write schematically
\begin{equation}
  \Sigma^{11,\mathrm{MR-}GW}
  =
  \Sigma_H^R
  +
  iG_DW_D
  \label{eq:MRGW-Hedin-form}
\end{equation}

Here
\begin{equation}
  W_D
  =
  v_R+w_D^c
\end{equation}
is the MR-RPA screened residual interaction.

The instantaneous bare part of $iG_DW_D$ generates the Fock exchange
contribution, while $iG_Dw_D^c$ generates the frequency-dependent
correlation self-energy.  Together with the Hartree term,
Eq.~\eqref{eq:MRGW-Hedin-form} is equivalent to the form
\begin{equation}
  \Sigma^{11,\mathrm{MR-}GW}
  =
  u
  +
  \bar v_R\gamma_D
  +
  iG_Dw_D^c
\end{equation}
used in the explicit MR-$GW$ expressions.


\subsubsection{Functional derivative of the parent self-energy}

The natural particle--hole irreducible kernel associated with the
physical self-energy is
\begin{equation}
  \mathcal I_D
  \equiv
  i
  \frac{
    \delta\Sigma^{11,\mathrm{MR-}GW}
  }{
    \delta G_D
  }
  \label{eq:kernel-functional-derivative}
\end{equation}
up to the conventional placement of factors of $i$ in the definition
of the Bethe--Salpeter kernel.

Differentiating Eq.~\eqref{eq:MRGW-Hedin-form} generates two classes of
terms:
\begin{equation}
  \frac{\delta\Sigma_H^R}{\delta G_D}
\end{equation}
and
\begin{equation}
  \frac{\delta(iG_DW_D)}{\delta G_D}
  =
  iW_D
  +
  iG_D
  \frac{\delta W_D}{\delta G_D}
\end{equation}

The first practical approximation is to neglect the variation of the
screened interaction,
\begin{equation}
  \frac{\delta W_D}{\delta G_D}
  \simeq
  0
  \label{eq:frozen-W}
\end{equation}

This is the usual ``frozen-$W$'' approximation entering the standard
$GW$--BSE kernel.

With the conventional particle--hole index routing, one obtains
\begin{equation}
  \begin{aligned}
  \mathcal I_D(1,2;3,4)
  \simeq{}&
  \delta(1,3)\delta(2,4)
  v_R(1,4)
  \\
  &-
  \delta(1,4)\delta(2,3)
  W_D(1,3)
  \end{aligned}
  \label{eq:MR-GW-BSE-kernel-dynamic}
\end{equation}

Equation~\eqref{eq:MR-GW-BSE-kernel-dynamic} is the direct
multireference analogue of the conventional $GW$--BSE kernel
\begin{equation}
  I=v-W
\end{equation}

The two terms have different origins: the first arises from the
functional derivative of the Hartree potential, while the second
arises from differentiating the explicit Green's-function line in the
$GW$ exchange-correlation self-energy.

They should therefore not be interpreted as two copies of the same
screened interaction.


\subsubsection{Static approximation}

A fully dynamical use of
Eq.~\eqref{eq:MR-GW-BSE-kernel-dynamic} would introduce a second time
difference into the mixed cumulant vertex.

The first-order mixed blocks contain the time-degenerate object
\begin{equation}
  \bar v_R C_2^D
\end{equation}
because the bare Coulomb interaction is instantaneous.  Replacing this
interaction by a fully dynamical $W_D(t,t')$ would instead require a
genuine non-degenerate three-time, or equivalently two-frequency,
Dyall vertex.

This would defeat the purpose of constructing the cheapest
beyond-MR-$GW$ approximation.

We therefore make the static approximation
\begin{equation}
  W_D
  \longrightarrow
  W_D^0
  \equiv
  W_D(\omega=0)
  \label{eq:static-WD}
\end{equation}

The resulting static kernel is
\begin{equation}
  \begin{aligned}
  \mathcal I_D^{\rm stat}(1,2;3,4)
  ={}&
  \delta(1,3)\delta(2,4)
  v_R(1,4)
  \\
  &-
  \delta(1,4)\delta(2,3)
  W_D^0(1,3)
  \end{aligned}
  \label{eq:static-MR-kernel}
\end{equation}

In orbital notation,
\begin{equation}
  [\mathcal I_D^{\rm stat}]_{pr,qs}
  =
  v^R_{pr,qs}
  -
  [W_D^0]_{ps,qr}
  \label{eq:static-MR-kernel-orbital}
\end{equation}

We will occasionally use the compact notation
\begin{equation}
  \mathcal I_D^{\rm stat}
  =
  v_R
  -
  (W_D^0)^x
  \label{eq:static-kernel-short}
\end{equation}
where the superscript $x$ denotes the crossed particle--hole index
pairing.

The term ``crossed'' is preferable here to ``direct'' or ``exchange'',
since the latter terminology depends on the representation and index
convention used for the four-point kernel.


\subsubsection{Relation to the bare first-order mixed block}

Using
\begin{equation}
  W_D^0
  =
  v_R
  +
  w_D^c(0)
\end{equation}
Eq.~\eqref{eq:static-kernel-short} may be written
\begin{equation}
  \mathcal I_D^{\rm stat}
  =
  \bar v_R
  -
  [w_D^c(0)]^x
  \label{eq:kernel-bare-plus-screening}
\end{equation}

Thus the proposed kernel differs from the exact first-order interaction
only through a static correlation correction in the crossed channel.

When screening is removed,
\begin{equation}
  W_D^0
  \longrightarrow
  v_R
\end{equation}
and therefore
\begin{equation}
  \mathcal I_D^{\rm stat}
  \longrightarrow
  v_R-v_R^x
  =
  \bar v_R
  \label{eq:kernel-bare-limit}
\end{equation}

Hence the first-order mixed self-energy derived in Sec.~\ref{sec:first-order-blocks} is recovered exactly.


\subsubsection{Screened mixed residues}

Let
\begin{equation}
  d_{\mu+;z}
  =
  \braket{
    \Xi_{\mu,N_a+1}^{A}
    |
    \hat z^\dagger
    |
    \Xi_{0,N_a}^{A}
  }
\end{equation}
and define the connected attachment tensor
\begin{equation}
  \mathcal C_{\mu;z,yw}^{D,+}
  =
  \frac12
  \braket{
    \Xi_{0,N_a}^{A}
    |
    \hat y^\dagger\hat w\hat z
    |
    \Xi_{\mu,N_a+1}^{A}
  }
  -
  \gamma_{yw}^{D}
  d_{\mu+;z}^{*}
  \label{eq:connected-attachment-tensor}
\end{equation}

Similarly, define the removal tensor
\begin{equation}
  \mathcal C_{\mu;z,yw}^{D,-}
  =
  \frac12
  \braket{
    \Xi_{\mu,N_a-1}^{A}
    |
    \hat y^\dagger\hat w\hat z
    |
    \Xi_{0,N_a}^{A}
  }
  -
  \gamma_{yw}^{D}
  d_{\mu-;z}
  \label{eq:connected-removal-tensor}
\end{equation}

The proposed statically screened mixed residues are then
\begin{equation}
  B^{+,{\rm scr}}_{P\mu}
  =
  [\mathcal I_D^{\rm stat}]_{Pz,yw}
  \mathcal C_{\mu;z,yw}^{D,+}
  \label{eq:Bplus-screened}
\end{equation}
and
\begin{equation}
  B^{-,{\rm scr}}_{P\mu}
  =
  [\mathcal I_D^{\rm stat}]_{Pz,yw}
  \mathcal C_{\mu;z,yw}^{D,-}
  \label{eq:Bminus-screened}
\end{equation}

For the explicit indices appearing in the first-order mixed blocks,
\begin{equation}
  [\mathcal I_D^{\rm stat}]_{Pz,yw}
  =
  v^R_{Pz,yw}
  -
  [W_D^0]_{Pw,yz}
  \label{eq:mixed-kernel-explicit}
\end{equation}

Equation~\eqref{eq:mixed-kernel-explicit} is therefore the practical
replacement
\begin{equation}
  \bar v^R_{Pz,yw}
  \quad\longrightarrow\quad
  v^R_{Pz,yw}
  -
  [W_D(0)]_{Pw,yz}
  \label{eq:screening-replacement}
\end{equation}
used in the first beyond-MR-$GW$ approximation.


\subsubsection{Why the fully screened antisymmetrized interaction is not
the default choice}

A seemingly natural alternative would be
\begin{equation}
  \mathcal I_D^{\rm full-W}
  =
  W_D^0-(W_D^0)^x
  \label{eq:full-W-alternative}
\end{equation}

This choice satisfies the same unscreened limit,
\begin{equation}
  \mathcal I_D^{\rm full-W}
  \longrightarrow
  \bar v_R
\end{equation}
and can therefore be considered as a useful numerical control.

It is, however, not the kernel obtained from the frozen-$W$ functional
derivative of the parent MR-$GW$ self-energy.

Indeed, the uncrossed term in
Eq.~\eqref{eq:static-MR-kernel-orbital} originates from the Hartree
functional derivative and is therefore the bare residual interaction,
whereas the crossed term originates from the $GW$ exchange-correlation
self-energy and contains $W_D$.

Replacing the first term by $W_D$ would therefore constitute an
additional approximation unrelated to the parent MR-$GW$ functional.

For this reason,
\begin{equation}
  \mathcal I_D^{\rm stat}
  =
  v_R-(W_D^0)^x
\end{equation}
is adopted as the default kernel.


\subsubsection{Relation to positive-semidefinite vertex constructions}

The same separation between a bare Coulomb interaction and a statically
screened interaction appears in positive-semidefinite extensions of
vertex-corrected $GW$.

There, the approximate four-point kernel has the structure
\begin{equation}
  I
  =
  v-W_0^x
\end{equation}
and the bare and statically screened interactions are treated as
distinct interactions when constructing the corresponding
half-amplitudes.

Our Eq.~\eqref{eq:static-kernel-short} is the direct Dyall-reference
analogue,
\begin{equation}
  v
  \rightarrow
  v_R
  \qquad
  W_0
  \rightarrow
  W_D(0)
\end{equation}

This analogy is useful but should not be confused with the logical
origin of the present kernel: the primary justification is the
functional derivative of the parent MR-$GW$ self-energy under the
frozen-$W$ and static approximations.


\subsubsection{Positive semidefiniteness does not determine the kernel}

It is important to separate two issues.

The choice
\begin{equation}
  \mathcal I_D^{\rm stat}
  =
  v_R-(W_D^0)^x
\end{equation}
determines the \emph{physical approximation to the mixed coupling}.

Positive semidefiniteness will instead be imposed at the next stage by
embedding these couplings into a Hermitian intermediate-state
representation.

Consequently, PSD alone does not select
Eq.~\eqref{eq:static-kernel-short}.  Other Hermitian static kernels
could also be given a PSD completion.

The present kernel is selected because it is simultaneously

\begin{itemize}
  \item consistent with the parent MR-$GW$ self-energy;
  \item obtained from the standard frozen-$W$ functional derivative;
  \item compatible with the conventional static $GW$--BSE
        approximation;
  \item exact in the bare first-order limit;
  \item inexpensive;
  \item and compatible with the subsequent Hermitian completion.
\end{itemize}


\subsubsection{Hierarchy of static mixed kernels}

For numerical analysis it is useful to distinguish three related
approximations.

\paragraph{Bare mixed kernel.}

\begin{equation}
  \mathcal I_D^{(0)}
  =
  \bar v_R
\end{equation}
This isolates the effect of the subsequent PSD/Hermitian completion
without introducing screening into the mixed block.

\paragraph{Parent-consistent statically screened kernel.}

\begin{equation}
  \mathcal I_D^{(1)}
  =
  v_R-(W_D^0)^x
  \label{eq:preferred-static-kernel}
\end{equation}
This is the default approximation.

\paragraph{Fully screened antisymmetrized kernel.}

\begin{equation}
  \mathcal I_D^{(2)}
  =
  W_D^0-(W_D^0)^x
\end{equation}
This is retained as a useful control but is not the functional
derivative of the parent MR-$GW$ approximation.


\subsubsection{Status of the approximation}

The choice
\begin{equation}
  \mathcal I_D^{\rm mix}
  =
  v_R-[W_D(0)]^x
\end{equation}
is not an exact identity of the Dyall-Hedin hierarchy.

It follows from the clearly defined sequence
\begin{equation}
  \begin{aligned}
  \text{Dyall-Hedin}
  &\longrightarrow
  \text{MR-}GW\text{ physical self-energy}
  \\
  &\longrightarrow
  \frac{\delta\Sigma^{\mathrm{MR-}GW}}{\delta G_D}
  \\
  &\xrightarrow{\delta W_D/\delta G_D\to0}
  v_R-W_D^x
  \\
  &\xrightarrow{W_D(\omega)\to W_D(0)}
  v_R-W_D^x(0)
  \end{aligned}
  \label{eq:kernel-derivation-summary}
\end{equation}

Within the frozen-screening and static approximations,
Eq.~\eqref{eq:preferred-static-kernel} is therefore the kernel
consistent with the parent MR-$GW$ theory.

Going beyond it requires restoring either
\begin{equation}
  \frac{\delta W_D}{\delta G_D}
\end{equation}
or the full frequency dependence of $W_D$, and therefore belongs to a
future dynamical MR-$GW\Gamma$ theory.
%%%%%%%%%%%%%%%%%%%%%%%%%
\subsection{Common Hermitian representation and PSD completion}
\label{sec:hermitian-completion}
%%%%%%%%%%%%%%%%%%%%%%%%%

\Proposed The finite Hermitian resolvent below defines our PSD
higher-order completion.  Several intermediate identities---in particular the
linearization of the parent MR-$GW$ self-energy and the first-order recovery of
the mixed Hall terms---are exact algebraic statements.


We now construct a common Hermitian representation containing both the
existing dynamic MR-$GW$ self-energy and the statically screened mixed
Hall blocks introduced in Sec.~\ref{sec:static-mixed-kernel}.

The purpose of this construction is threefold:
\begin{enumerate}
  \item to recover the existing MR-$GW$ approximation exactly when the
  mixed couplings are switched off;
  \item to reproduce the screened $\Sigma^{12}$ and $\Sigma^{21}$
  contributions through first order;
  \item to define a higher-order completion with a positive-semidefinite
  spectral function and the correct zeroth spectral moment.
\end{enumerate}

The construction is most transparent for the retarded Green's function.
Throughout this section,
\begin{equation}
  z
  =
  \omega+i0^+
\end{equation}


\subsubsection{Retarded spectral representation of the Dyall Green's function}

Let $n$ denote collectively all poles of the exact Dyall Green's
function.  These include
\begin{itemize}
  \item inactive virtual poles;
  \item inactive core-removal poles;
  \item active-space $N+1$ poles;
  \item active-space $N-1$ poles.
\end{itemize}

We introduce a signed pole energy $\kappa_n$ for every pole and collect
these energies in the Hermitian matrix
\begin{equation}
  K_D
  =
  \operatorname{diag}
  \left(
    \kappa_1,\kappa_2,\ldots
  \right)
\end{equation}

Let $t_n$ denote the corresponding physical transition vector.  The
retarded Dyall Green's function can then be written
\begin{equation}
  G_D^R(z)
  =
  \sum_n
  \frac{
    t_nt_n^\dagger
  }{
    z-\kappa_n
  }
  \label{eq:GD-retarded-spectral}
\end{equation}

Introducing the transition matrix
\begin{equation}
  T_D
  =
  \begin{pmatrix}
    t_1&t_2&\cdots
  \end{pmatrix}
\end{equation}
and the pole-space resolvent
\begin{equation}
  R_D(z)
  =
  \left(
    zI-K_D
  \right)^{-1}
\end{equation}
Eq.~\eqref{eq:GD-retarded-spectral} becomes
\begin{equation}
  G_D^R(z)
  =
  T_DR_D(z)T_D^\dagger
  \label{eq:GD-TKT}
\end{equation}

Completeness of the exact $N+1$ and $N-1$ Dyall states gives the
fermionic zeroth-moment sum rule
\begin{equation}
  T_DT_D^\dagger
  =
  I
  \label{eq:TD-isometry}
\end{equation}

Importantly, $T_D$ need not be square.  In a correlated reference the
number of Green's-function poles generally exceeds the dimension of
the physical one-particle space, so that in general
\begin{equation}
  T_D^\dagger T_D
  \neq
  I
\end{equation}


\subsubsection{Inactive and active pole states}

For an inactive orbital $P$, there is a single zeroth-order pole with
transition vector
\begin{equation}
  t_P=e_P
\end{equation}
where $e_P$ is the unit vector associated with orbital $P$, and
\begin{equation}
  \kappa_P=\epsilon_P
\end{equation}

For an active charged state we introduce the composite index
\begin{equation}
  \alpha
  \equiv
  (\mu,\sigma)
  \qquad
  \sigma=\pm
\end{equation}
with transition vector
\begin{equation}
  d_\alpha
  \equiv
  t_{\mu\sigma}
\end{equation}

The corresponding signed pole energies are
\begin{equation}
  \kappa_{\mu+}
  =
  E_{\mu}^{N+1}-E_0^N
\end{equation}
and
\begin{equation}
  \kappa_{\mu-}
  =
  -
  \left(
    E_{\mu}^{N-1}-E_0^N
  \right)
\end{equation}

Hence
\begin{equation}
  G_D^{A,R}(z)
  =
  \sum_\alpha
  \frac{
    d_\alpha d_\alpha^\dagger
  }{
    z-\kappa_\alpha
  }
\end{equation}


\subsubsection{A useful linearization identity}

Before introducing the mixed blocks, it is useful to establish a
general identity.

Let
\begin{equation}
  G_0(z)
  =
  TR(z)T^\dagger
\end{equation}
with
\begin{equation}
  TT^\dagger=I
\end{equation}
and let $X(z)$ be an arbitrary operator in the physical one-particle
space.

Then
\begin{equation}
  T
  \left[
    R^{-1}(z)
    -
    T^\dagger X(z)T
  \right]^{-1}
  T^\dagger
  =
  \left[
    G_0^{-1}(z)-X(z)
  \right]^{-1}
  \label{eq:linearization-lemma}
\end{equation}

To prove Eq.~\eqref{eq:linearization-lemma}, use the Woodbury identity,
\begin{equation}
  \begin{aligned}
  &
  \left[
    R^{-1}
    -
    T^\dagger XT
  \right]^{-1}
  \\
  &\qquad
  =
  R
  +
  RT^\dagger
  X
  \left(
    I-TRT^\dagger X
  \right)^{-1}
  TR
  \end{aligned}
\end{equation}
Multiplication from the left and right by $T$ and $T^\dagger$ gives
\begin{equation}
  \begin{aligned}
  T
  \left[
    R^{-1}
    -
    T^\dagger XT
  \right]^{-1}
  T^\dagger
  ={}&
  G_0
  +
  G_0X
  \left(
    I-G_0X
  \right)^{-1}
  G_0
  \\
  ={}&
  \left(
    I-G_0X
  \right)^{-1}
  G_0
  \\
  ={}&
  \left(
    G_0^{-1}-X
  \right)^{-1}
  \end{aligned}
\end{equation}

Equation~\eqref{eq:linearization-lemma} remains valid when $T$ is
rectangular.  This is crucial for an interacting reference with more
poles than physical orbitals.


\subsubsection{Pole representation of the MR-$GW$ self-energy}

The retarded correlation self-energy of the existing MR-$GW$
approximation has a simple-pole representation,
\begin{equation}
  \Sigma_{\rm MR-GW}^{c,R}(z)
  =
  A
  \left(
    zI-E_{\rm GW}
  \right)^{-1}
  A^\dagger
  \label{eq:MRGW-self-energy-pole}
\end{equation}
where $E_{\rm GW}$ is a real diagonal matrix containing the signed
energies of the MR-$GW$ intermediate states.

This follows directly from the spectral expression derived in
Ref.~\cite{WangFangLi2026}.

For example, the pole energies include
\begin{subequations}
  \begin{align}
    E_{aI}
    &=
    \epsilon_a+\Omega_I
    \\
    E_{\mu+I}
    &=
    \omega_{\mu}^{N+1}+\Omega_I
    \\
    E_{iI}
    &=
    \epsilon_i-\Omega_I
    \\
    E_{\mu-I}
    &=
    -
    \left(
      \omega_{\mu}^{N-1}+\Omega_I
    \right)
  \end{align}
\end{subequations}

For real orbitals and amplitudes, representative coupling amplitudes
may be chosen as
\begin{subequations}
  \begin{align}
    A_{p,aI}
    &=
    M_{ap,I}
    \\
    A_{p,\mu+I}
    &=
    D_{\mu,x}^{[+1]}
    M_{xp,I}
    \\
    A_{p,iI}
    &=
    M_{pi,I}
    \\
    A_{p,\mu-I}
    &=
    D_{\mu,x}^{[-1]}
    M_{px,I}
  \end{align}
\end{subequations}
where repeated active indices are summed.

For complex quantities these definitions are understood as any
factorization satisfying
\begin{equation}
  \operatorname*{Res}_{z=E_\Lambda}
  \Sigma_{\rm MR-GW}^{c,R}(z)
  =
  A_\Lambda A_\Lambda^\dagger
\end{equation}

For a stable MR-RPA solution these residue matrices are positive
semidefinite.

We separate the static contribution,
\begin{equation}
  \Sigma_{\rm MR-GW}^{R}(z)
  =
  \Sigma_{\rm stat}
  +
  A
  \left(
    zI-E_{\rm GW}
  \right)^{-1}
  A^\dagger
  \label{eq:full-MRGW-selfenergy-pole}
\end{equation}
where
\begin{equation}
  \Sigma_{\rm stat}
  \equiv
  \Sigma_1^{11}
\end{equation}


\subsubsection{Exact Hermitian linearization of the existing MR-$GW$ approximation}

Define the enlarged Hermitian matrix
\begin{equation}
  \mathbb H_{\rm MR-GW}
  =
  \begin{pmatrix}
    K_D
    +
    T_D^\dagger
    \Sigma_{\rm stat}
    T_D
    &
    T_D^\dagger A
    \\[2mm]
    A^\dagger T_D
    &
    E_{\rm GW}
  \end{pmatrix}
  \label{eq:H-MRGW}
\end{equation}

The physical transition matrix into the enlarged space is
\begin{equation}
  \mathbb T
  =
  \begin{pmatrix}
    T_D & 0
  \end{pmatrix}
  \label{eq:T-enlarged}
\end{equation}

Consider
\begin{equation}
  G_{\rm lin}^{R}(z)
  =
  \mathbb T
  \left(
    zI-\mathbb H_{\rm MR-GW}
  \right)^{-1}
  \mathbb T^\dagger
  \label{eq:G-linearized-MRGW}
\end{equation}

Eliminating the MR-$GW$ bath by a Schur complement gives
\begin{equation}
  \begin{aligned}
  G_{\rm lin}^{R}(z)
  =
  T_D
  \Big[
  &
  R_D^{-1}(z)
  -
  T_D^\dagger
  \Sigma_{\rm stat}
  T_D
  \\
  &
  -
  T_D^\dagger
  A
  \left(
    zI-E_{\rm GW}
  \right)^{-1}
  A^\dagger
  T_D
  \Big]^{-1}
  T_D^\dagger
  \end{aligned}
\end{equation}

Using Eq.~\eqref{eq:linearization-lemma} and
Eq.~\eqref{eq:full-MRGW-selfenergy-pole},
\begin{equation}
  G_{\rm lin}^{R}(z)
  =
  \left[
    (G_D^R)^{-1}(z)
    -
    \Sigma_{\rm MR-GW}^{R}(z)
  \right]^{-1}
  \label{eq:linearization-equals-Dyson}
\end{equation}

Thus
\begin{equation}
  \mathbb H_{\rm MR-GW}
\end{equation}
is an exact Hermitian linearization of the published MR-$GW$ Dyson
equation.

No approximation has been introduced in passing from the pole
representation of the MR-$GW$ self-energy to
Eq.~\eqref{eq:H-MRGW}.


\subsubsection{Retarded representation of the screened mixed blocks}

The statically screened mixed blocks derived in the previous section
have the same active-space pole structure as the first-order blocks.

After conversion to retarded form, write
\begin{equation}
  [\Sigma_{\rm scr}^{12,R}]_{Px}(z)
  =
  \sum_\alpha
  \frac{
    B_{P\alpha}^{\rm scr}
    [d_\alpha^\dagger]_x
  }{
    z-\kappa_\alpha
  }
  \label{eq:sigma12-retarded-B}
\end{equation}

The corresponding left block is
\begin{equation}
  [\Sigma_{\rm scr}^{21,R}]_{xP}(z)
  =
  \sum_\alpha
  \frac{
    [d_\alpha]_x
    B_{P\alpha}^{{\rm scr}*}
  }{
    z-\kappa_\alpha
  }
  \label{eq:sigma21-retarded-B}
\end{equation}

The index $\alpha$ runs over \emph{both} active addition and active
removal poles.

The screened residue is
\begin{equation}
  B_{P\alpha}^{\rm scr}
  =
  \mathcal I_D^{\rm stat}
  \cdot
  \mathcal C_{P\alpha}^{D}
\end{equation}
with
\begin{equation}
  \mathcal I_D^{\rm stat}
  =
  v_R-[W_D(0)]^x
\end{equation}


\subsubsection{Hermitian pole-space representation of the mixed blocks}

Define a matrix $K_B$ in the Dyall pole space by
\begin{equation}
  [K_B]_{P\alpha}
  =
  B_{P\alpha}^{\rm scr}
  \label{eq:KB-Palpha}
\end{equation}
and
\begin{equation}
  [K_B]_{\alpha P}
  =
  B_{P\alpha}^{{\rm scr}*}
  \label{eq:KB-alphaP}
\end{equation}

All other matrix elements are set to zero in the minimal approximation.

Therefore
\begin{equation}
  K_B^\dagger=K_B
  \label{eq:KB-Hermitian}
\end{equation}

We now evaluate the first-order correction generated by $K_B$,
\begin{equation}
  \delta G_B^R(z)
  =
  T_DR_D(z)
  K_B
  R_D(z)T_D^\dagger
  \label{eq:dG-KB}
\end{equation}

Using the explicit inactive and active pole vectors,
Eq.~\eqref{eq:dG-KB} gives
\begin{equation}
  \begin{aligned}
  \delta G_B^R(z)
  ={}&
  \sum_{P\alpha}
  \frac{
    e_P
    B_{P\alpha}^{\rm scr}
    d_\alpha^\dagger
  }{
    (z-\kappa_P)
    (z-\kappa_\alpha)
  }
  \\
  &+
  \sum_{P\alpha}
  \frac{
    d_\alpha
    B_{P\alpha}^{{\rm scr}*}
    e_P^\dagger
  }{
    (z-\kappa_\alpha)
    (z-\kappa_P)
  }
  \end{aligned}
  \label{eq:dG-KB-explicit}
\end{equation}


\subsubsection{Exact recovery of the mixed first-order Green's-function terms}

From Eq.~\eqref{eq:sigma12-retarded-B},
\begin{equation}
  \begin{aligned}
  [G_D^R\Sigma_{\rm scr}^{12,R}]_{Px}
  &=
  \frac{1}{z-\kappa_P}
  [\Sigma_{\rm scr}^{12,R}]_{Px}
  \\
  &=
  \sum_\alpha
  \frac{
    B_{P\alpha}^{\rm scr}
    [d_\alpha^\dagger]_x
  }{
    (z-\kappa_P)
    (z-\kappa_\alpha)
  }
  \end{aligned}
\end{equation}

Likewise,
\begin{equation}
  \begin{aligned}
  [\Sigma_{\rm scr}^{21,R}G_D^R]_{xP}
  &=
  [\Sigma_{\rm scr}^{21,R}]_{xP}
  \frac{1}{z-\kappa_P}
  \\
  &=
  \sum_\alpha
  \frac{
    [d_\alpha]_x
    B_{P\alpha}^{{\rm scr}*}
  }{
    (z-\kappa_\alpha)
    (z-\kappa_P)
  }
  \end{aligned}
\end{equation}

Comparing with Eq.~\eqref{eq:dG-KB-explicit} gives the exact identity
\begin{equation}
  T_DR_DK_BR_DT_D^\dagger
  =
  G_D^R\Sigma_{\rm scr}^{12,R}
  +
  \Sigma_{\rm scr}^{21,R}G_D^R
  \label{eq:KB-equals-mixed}
\end{equation}

Equation~\eqref{eq:KB-equals-mixed} is the central algebraic result of
the present section.

A single Hermitian pole-space matrix $K_B$ reproduces both mixed Hall
blocks through first order.


\subsubsection{Same-sector pole pairs}

Suppose that the inactive pole $P$ and the active pole $\alpha$ belong
to the same physical charged sector.

For example,
\begin{equation}
  P=a
  \qquad
  \alpha=(\mu,+)
\end{equation}
in the electron-addition sector, or
\begin{equation}
  P=i
  \qquad
  \alpha=(\mu,-)
\end{equation}
in the electron-removal sector.

Then
\begin{equation}
  [K_B]_{P\alpha}
\end{equation}
has the natural interpretation of a screened coupling between Dyall
states belonging to the same $N\pm1$ sector.

In the bare limit,
\begin{equation}
  W_D(0)\rightarrow v_R
\end{equation}
the combination of this mixed residue with the corresponding
off-diagonal contribution from $\Sigma_1^{11}$ reconstructs the
first-order residual-Hamiltonian coupling between the two charged
Dyall states.

Thus the Hermitian representation has a direct charged-state
interpretation for same-sector pole pairs.


\subsubsection{Cross-sector pole pairs}

The situation is slightly different when the two poles belong to
opposite Lehmann sectors, for example
\begin{equation}
  P=i
  \qquad
  \alpha=(\mu,+)
\end{equation}
or
\begin{equation}
  P=a
  \qquad
  \alpha=(\mu,-)
\end{equation}

Such poles cannot be interpreted as physical states belonging to the
same fixed-particle-number Hamiltonian.

Nevertheless, Eq.~\eqref{eq:KB-equals-mixed} remains valid without
modification.

The reason is that the Hermitian matrix introduced here is an
\emph{auxiliary pole-space representation of the retarded Green's
function}; it is not the physical $N+1$ or $N-1$ Hamiltonian.

All retarded poles occur with the same analytic prescription
\begin{equation}
  \frac{1}{z-\kappa_n}
  \qquad
  z=\omega+i0^+
\end{equation}
so a Hermitian auxiliary coupling may connect any pair of signed poles.


\subsubsection{Equivalent first-order transition-amplitude representation}

The cross-sector structure may also be understood by partial
fractioning.

For $\kappa_P\neq\kappa_\alpha$,
\begin{equation}
  \frac{1}{
    (z-\kappa_P)
    (z-\kappa_\alpha)
  }
  =
  \frac{1}{
    \kappa_P-\kappa_\alpha
  }
  \left[
    \frac{1}{z-\kappa_P}
    -
    \frac{1}{z-\kappa_\alpha}
  \right]
  \label{eq:cross-partial-fraction}
\end{equation}

Define
\begin{equation}
  \Delta_{P\alpha}
  =
  \kappa_P-\kappa_\alpha
\end{equation}

For each cross-sector pair, introduce
\begin{equation}
  S_{\alpha P}
  =
  \frac{
    B_{P\alpha}^{{\rm scr}*}
  }{
    \Delta_{P\alpha}
  }
  \label{eq:S-alphaP}
\end{equation}
and
\begin{equation}
  S_{P\alpha}
  =
  -
  \frac{
    B_{P\alpha}^{\rm scr}
  }{
    \Delta_{P\alpha}
  }
  \label{eq:S-Palpha}
\end{equation}

Since the pole energies are real,
\begin{equation}
  S^\dagger=-S
  \label{eq:S-antihermitian}
\end{equation}

A first-order rotation of the transition matrix,
\begin{equation}
  \delta T_D
  =
  T_DS
  \label{eq:dT-TS}
\end{equation}
therefore gives
\begin{equation}
  \delta t_P
  =
  d_\alpha
  \frac{
    B_{P\alpha}^{{\rm scr}*}
  }{
    \Delta_{P\alpha}
  }
\end{equation}
and
\begin{equation}
  \delta d_\alpha
  =
  -
  e_P
  \frac{
    B_{P\alpha}^{\rm scr}
  }{
    \Delta_{P\alpha}
  }
\end{equation}

The corresponding residue corrections at $\kappa_P$ and
$\kappa_\alpha$ reproduce exactly the partial-fraction form of
Eq.~\eqref{eq:dG-KB-explicit}.

Thus the two representations
\begin{equation}
  K_B
  \quad\text{and}\quad
  \delta T_D=T_DS
\end{equation}
are equivalent through first order for cross-sector contributions.


\subsubsection{Preservation of the zeroth moment in the transition-rotation representation}

Because $S$ is anti-Hermitian, the finite transformation
\begin{equation}
  \widetilde T_D
  =
  T_De^S
  \label{eq:T-unitary-rotation}
\end{equation}
is unitary in pole space.

Hence
\begin{equation}
  \widetilde T_D
  \widetilde T_D^\dagger
  =
  T_De^Se^{-S}T_D^\dagger
  =
  T_DT_D^\dagger
  =
  I
\end{equation}

Therefore the cross-sector transition-amplitude completion preserves
the zeroth spectral-moment sum rule exactly.

This provides an independent check on the sign and denominator in
Eqs.~\eqref{eq:S-alphaP} and \eqref{eq:S-Palpha}.


\subsubsection{Choice of higher-order completion}

Although the Hamiltonian-coupling representation and the
transition-rotation representation are equivalent through first order,
their finite higher-order completions are not identical.

The present method adopts the common Hermitian pole-space coupling
$K_B$ because it
\begin{itemize}
  \item treats same- and cross-sector pole pairs uniformly;
  \item combines naturally with the Hermitian MR-$GW$ pole bath;
  \item preserves causality and positive spectral weight;
  \item and avoids a separate cross-sector implementation.
\end{itemize}

The resulting higher-order terms should therefore be regarded as a
\emph{chosen Hermitian resummation} of the screened mixed blocks, not
as the unique higher-order continuation implied by the exact
Dyall-Hedin hierarchy.


\subsubsection{Common Hermitian matrix}

We can now combine the existing MR-$GW$ approximation with the
screened mixed couplings.

Define
\begin{equation}
  K_{\rm top}
  =
  K_D
  +
  T_D^\dagger
  \Sigma_{\rm stat}
  T_D
  +
  K_B
  \label{eq:K-top-final}
\end{equation}

The complete auxiliary Hamiltonian is
\begin{equation}
  \mathbb H_{\rm mix-MR-GW}
  =
  \begin{pmatrix}
    K_{\rm top}
    &
    T_D^\dagger A
    \\[2mm]
    A^\dagger T_D
    &
    E_{\rm GW}
  \end{pmatrix}
  \label{eq:H-common-final}
\end{equation}

By construction,
\begin{equation}
  \mathbb H_{\rm mix-MR-GW}^\dagger
  =
  \mathbb H_{\rm mix-MR-GW}
  \label{eq:H-common-Hermitian}
\end{equation}

The physical Green's function defining the proposed approximation is
\begin{equation}
  G_{\rm mix-MR-GW}^{R}(z)
  =
  \mathbb T
  \left(
    zI-\mathbb H_{\rm mix-MR-GW}
  \right)^{-1}
  \mathbb T^\dagger
  \label{eq:G-common-final}
\end{equation}
with
\begin{equation}
  \mathbb T
  =
  \begin{pmatrix}
    T_D&0
  \end{pmatrix}
\end{equation}


\subsubsection{Recovery of the existing MR-$GW$ approximation}

If the mixed coupling is switched off,
\begin{equation}
  K_B=0
\end{equation}
then
\begin{equation}
  \mathbb H_{\rm mix-MR-GW}
  \longrightarrow
  \mathbb H_{\rm MR-GW}
\end{equation}

Using the exact linearization derived above,
\begin{equation}
  G_{\rm mix-MR-GW}^{R}(z)
  \xrightarrow{K_B\to0}
  \left[
    (G_D^R)^{-1}(z)
    -
    \Sigma_{\rm MR-GW}^{R}(z)
  \right]^{-1}
  \label{eq:KB-zero-MRGW}
\end{equation}

Thus the existing MR-$GW$ approximation is recovered exactly, not only
perturbatively.


\subsubsection{First-order expansion of the common Hermitian representation}

To verify the perturbative limit, introduce a formal parameter
$\lambda$ multiplying the residual interaction.

The relevant quantities scale as
\begin{equation}
  \Sigma_{\rm stat}
  =
  \mathcal O(\lambda)
\end{equation}
\begin{equation}
  K_B
  =
  \mathcal O(\lambda)
  +
  \text{screening corrections}
\end{equation}
and
\begin{equation}
  A
  =
  \mathcal O(\lambda)
\end{equation}
so that
\begin{equation}
  A
  \left(
    zI-E_{\rm GW}
  \right)^{-1}
  A^\dagger
  =
  \mathcal O(\lambda^2)
\end{equation}

Expansion of Eq.~\eqref{eq:G-common-final} therefore gives
\begin{equation}
  \begin{aligned}
  G_{\rm mix-MR-GW}^{R}
  =
  G_D^R
  &+
  G_D^R
  \Sigma_{\rm stat}
  G_D^R
  \\
  &+
  T_DR_DK_BR_DT_D^\dagger
  +
  \mathcal O(\lambda^2)
  \end{aligned}
\end{equation}

Using Eq.~\eqref{eq:KB-equals-mixed},
\begin{equation}
  \begin{aligned}
  G_{\rm mix-MR-GW}^{R}
  =
  G_D^R
  &+
  G_D^R
  \Sigma_{\rm stat}
  G_D^R
  \\
  &+
  G_D^R
  \Sigma_{\rm scr}^{12,R}
  +
  \Sigma_{\rm scr}^{21,R}
  G_D^R
  +
  \mathcal O(\lambda^2)
  \end{aligned}
  \label{eq:firstorder-common}
\end{equation}

In the unscreened limit,
\begin{equation}
  W_D(0)\rightarrow v_R
\end{equation}
the screened mixed residues reduce to the exact first-order residues,
\begin{equation}
  \Sigma_{\rm scr}^{12}
  \rightarrow
  \Sigma_1^{12}
\end{equation}
and
\begin{equation}
  \Sigma_{\rm scr}^{21}
  \rightarrow
  \Sigma_1^{21}
\end{equation}

Since
\begin{equation}
  \Sigma_1^{22}=0
\end{equation}
for the Dyall Hamiltonian, Eq.~\eqref{eq:firstorder-common} becomes
\begin{equation}
  G_1
  =
  G_D\Sigma_1^{11}G_D
  +
  G_D\Sigma_1^{12}
  +
  \Sigma_1^{21}G_D
\end{equation}
which is the complete exact first-order Green's function.

Thus the Hermitian mixed-block extension restores the first-order
information discarded by the practical MR-$GW$ approximation.


\subsubsection{Positive-semidefinite spectral function}

Since
\begin{equation}
  \mathbb H_{\rm mix-MR-GW}
  =
  \mathbb H_{\rm mix-MR-GW}^\dagger
\end{equation}
it admits an orthonormal eigenbasis,
\begin{equation}
  \mathbb H_{\rm mix-MR-GW}
  U_\lambda
  =
  E_\lambda U_\lambda
\end{equation}

Equation~\eqref{eq:G-common-final} becomes
\begin{equation}
  G_{\rm mix-MR-GW}^{R}(z)
  =
  \sum_\lambda
  \frac{
    Z_\lambda Z_\lambda^\dagger
  }{
    z-E_\lambda
  }
  \label{eq:G-positive-Lehmann}
\end{equation}
where
\begin{equation}
  Z_\lambda
  =
  \mathbb T U_\lambda
\end{equation}

Every residue matrix therefore satisfies
\begin{equation}
  Z_\lambda Z_\lambda^\dagger
  \succeq
  0
\end{equation}

Consequently,
\begin{equation}
  A(\omega)
  =
  -\frac{1}{\pi}
  \operatorname{Im}
  G^R(\omega)
  \succeq
  0
  \label{eq:spectral-PSD}
\end{equation}

The proposed Hermitian completion therefore possesses a
positive-semidefinite spectral function by construction.


\subsubsection{Zeroth spectral moment}

Completeness of the eigenvectors of
$\mathbb H_{\rm mix-MR-GW}$ gives
\begin{equation}
  \sum_\lambda
  U_\lambda U_\lambda^\dagger
  =
  I
\end{equation}

Hence
\begin{equation}
  \begin{aligned}
  \sum_\lambda
  Z_\lambda Z_\lambda^\dagger
  &=
  \mathbb T
  \left(
    \sum_\lambda
    U_\lambda U_\lambda^\dagger
  \right)
  \mathbb T^\dagger
  \\
  &=
  \mathbb T\mathbb T^\dagger
  \\
  &=
  T_DT_D^\dagger
  \\
  &=
  I
  \end{aligned}
\end{equation}

Therefore
\begin{equation}
  \int_{-\infty}^{+\infty}
  d\omega\
  A(\omega)
  =
  I
  \label{eq:zeroth-moment-exact}
\end{equation}

Thus the Hermitian completion preserves both positive spectral weight
and the exact fermionic zeroth-moment sum rule.


\subsubsection{Causality}

Equation~\eqref{eq:G-common-final} is the resolvent of a Hermitian
matrix evaluated at
\begin{equation}
  z=\omega+i0^+
\end{equation}

It therefore has no poles in the upper half of the complex frequency
plane and obeys the usual retarded causality condition.

Hence the Hermitian completion simultaneously guarantees
\begin{equation}
  \text{causality}
  +
  \text{PSD}
  +
  \text{zeroth-moment normalization}
\end{equation}


\subsubsection{Scope of the result}

The following statements are exact algebraic results:

\begin{enumerate}
  \item the matrix $\mathbb H_{\rm MR-GW}$ is an exact linearization of
  the published MR-$GW$ Dyson equation;
  \item $K_B$ reproduces
  \begin{equation}
    G_D\Sigma_{\rm scr}^{12}
    +
    \Sigma_{\rm scr}^{21}G_D
  \end{equation}
  exactly through first order;
  \item this statement holds for both same-sector and cross-sector
  pole pairs;
  \item the resulting common resolvent is causal, PSD, and preserves
  the zeroth spectral moment.
\end{enumerate}

What is \emph{not} uniquely determined by the exact first-order theory
is the higher-order resummation generated by repeated insertions of
$K_B$.

The choice
\begin{equation}
  \left(
    zI-\mathbb H_{\rm mix-MR-GW}
  \right)^{-1}
\end{equation}
therefore defines a specific Hermitian/PSD higher-order completion of
the screened mixed blocks.

This completion is selected because it preserves all known first-order
constraints while avoiding the pathological incomplete resummation
produced by inserting isolated first-order Hall blocks into the
nonlinear generalized Dyson equation.
%%%%%%%%%%%%%%%%%%%%%%%%%
\subsection{Spectral moments}
\label{sec:first-spectral-moment}
%%%%%%%%%%%%%%%%%%%%%%%%%

\Exact The spectral-moment identities themselves are exact.  Their
evaluation with the proposed Hermitian Hamiltonian quantifies which moment
constraints are preserved by the approximation.


The Hermitian representation introduced in
Sec.~\ref{sec:hermitian-completion} guarantees both positive spectral
weight and the exact zeroth spectral moment.  We now examine the first
spectral moment,
\begin{equation}
  M_1
  =
  \int_{-\infty}^{+\infty}
  d\omega\
  \omega A(\omega)
  \label{eq:first-moment-def}
\end{equation}
which provides a more stringent diagnostic of the approximation.

For the exact retarded Green's function, the first spectral moment is
determined by the equal-time equation of motion,
\begin{equation}
  [M_1^{\rm exact}]_{pq}
  =
  \braket{
    \left\{
      [\hat a_p,\hat H]
      \hat a_q^\dagger
    \right\}
  }
  \label{eq:exact-first-moment-commutator}
\end{equation}
Consequently, the first moment probes the high-frequency static
structure of the exact one-particle Green's function.


\subsubsection{First moment from the high-frequency expansion}

Consider a Green's function written in the Hermitian form
\begin{equation}
  G^R(z)
  =
  \mathbb T
  \left(
    zI-\mathbb H
  \right)^{-1}
  \mathbb T^\dagger
  \qquad
  z=\omega+i0^+
  \label{eq:G-Hermitian-general-moment}
\end{equation}

For large $|z|$,
\begin{equation}
  \left(
    zI-\mathbb H
  \right)^{-1}
  =
  \frac{I}{z}
  +
  \frac{\mathbb H}{z^2}
  +
  \frac{\mathbb H^2}{z^3}
  +
  \cdots
\end{equation}

Therefore,
\begin{equation}
  G^R(z)
  =
  \frac{
    \mathbb T\mathbb T^\dagger
  }{z}
  +
  \frac{
    \mathbb T\mathbb H\mathbb T^\dagger
  }{z^2}
  +
  \mathcal O(z^{-3})
  \label{eq:G-high-frequency}
\end{equation}

Comparison with the spectral expansion
\begin{equation}
  G^R(z)
  =
  \frac{M_0}{z}
  +
  \frac{M_1}{z^2}
  +
  \frac{M_2}{z^3}
  +
  \cdots
\end{equation}
gives
\begin{equation}
  M_0
  =
  \mathbb T\mathbb T^\dagger
\end{equation}
and
\begin{equation}
  M_1
  =
  \mathbb T
  \mathbb H
  \mathbb T^\dagger
  \label{eq:M1-Hermitian}
\end{equation}


\subsubsection{First moment of the exact Dyall reference}

For the Dyall reference,
\begin{equation}
  G_D^R(z)
  =
  T_D
  \left(
    zI-K_D
  \right)^{-1}
  T_D^\dagger
\end{equation}

Its first moment is therefore
\begin{equation}
  M_1^D
  =
  T_DK_DT_D^\dagger
  \label{eq:M1-Dyall}
\end{equation}

Since $G_D$ is the exact Green's function of $\hat H_D$,
Eq.~\eqref{eq:M1-Dyall} is equivalently
\begin{equation}
  [M_1^D]_{pq}
  =
  \braket{
    \left\{
      [\hat a_p,\hat H_D]
      \hat a_q^\dagger
    \right\}
  }_D
  \label{eq:M1-Dyall-commutator}
\end{equation}


\subsubsection{Exact first moment for a two-body Hamiltonian}

For reference, consider a Hamiltonian written as
\begin{equation}
  \hat H
  =
  h_{pq}
  \hat a_p^\dagger\hat a_q
  +
  \frac14
  \bar v_{pq,rs}
  \hat a_p^\dagger
  \hat a_q^\dagger
  \hat a_s
  \hat a_r
  \label{eq:H-two-body-moment}
\end{equation}

Using Eq.~\eqref{eq:exact-first-moment-commutator}, one obtains
\begin{equation}
  [M_1^{\rm exact}]_{ij}
  =
  h_{ij}
  +
  \bar v_{ir,js}
  \gamma_{rs}^{\rm exact}
  \label{eq:M1-exact-two-body}
\end{equation}
where
\begin{equation}
  \gamma_{rs}^{\rm exact}
  =
  \braket{
    \hat a_r^\dagger
    \hat a_s
  }
\end{equation}

Thus the exact first moment depends only on the exact one-body density
matrix, despite the presence of a two-body interaction.

Equivalently, if the exact self-energy has the high-frequency
expansion
\begin{equation}
  \Sigma(z)
  =
  \Sigma_\infty
  +
  \frac{\Sigma_1}{z}
  +
  \cdots
\end{equation}
then
\begin{equation}
  \Sigma_\infty
\end{equation}
is precisely the Hartree--Fock-like contraction of the bare
interaction with the exact one-body density matrix.


\subsubsection{First moment of the proposed Hermitian approximation}

The common Hermitian matrix obtained in Sec.~\ref{sec:hermitian-completion} is
\begin{equation}
  \mathbb H_{\rm mix-MR-GW}
  =
  \begin{pmatrix}
    K_{\rm top}
    &
    T_D^\dagger A
    \\[2mm]
    A^\dagger T_D
    &
    E_{\rm GW}
  \end{pmatrix}
  \label{eq:H-task6}
\end{equation}
with
\begin{equation}
  K_{\rm top}
  =
  K_D
  +
  T_D^\dagger
  \Sigma_{\rm stat}
  T_D
  +
  K_B
  \label{eq:Ktop-task6}
\end{equation}
and
\begin{equation}
  \mathbb T
  =
  \begin{pmatrix}
    T_D&0
  \end{pmatrix}
\end{equation}

Using Eq.~\eqref{eq:M1-Hermitian},
\begin{equation}
  M_1^{\rm mix-MR-GW}
  =
  T_DK_{\rm top}T_D^\dagger
\end{equation}

Substitution of Eq.~\eqref{eq:Ktop-task6} yields
\begin{equation}
  \begin{aligned}
  M_1^{\rm mix-MR-GW}
  ={}&
  T_DK_DT_D^\dagger
  \\
  &+
  T_DT_D^\dagger
  \Sigma_{\rm stat}
  T_DT_D^\dagger
  \\
  &+
  T_DK_BT_D^\dagger
  \end{aligned}
\end{equation}

Using the exact zeroth-moment identity
\begin{equation}
  T_DT_D^\dagger=I
\end{equation}
we obtain the central result
\begin{equation}
  M_1^{\rm mix-MR-GW}
  =
  M_1^D
  +
  \Sigma_{\rm stat}
  +
  T_DK_BT_D^\dagger
  \label{eq:M1-main-result}
\end{equation}


\subsubsection{The dynamic MR-$GW$ bath does not affect the first moment}

A notable consequence of Eq.~\eqref{eq:M1-main-result} is the complete
absence of
\begin{equation}
  A
  \qquad\text{and}\qquad
  E_{\rm GW}
\end{equation}

The dynamic MR-$GW$ correlation self-energy therefore does not modify
the first spectral moment.

This can also be seen directly from its high-frequency expansion,
\begin{equation}
  \begin{aligned}
  \Sigma_{\rm MR-GW}^{c}(z)
  &=
  A
  \left(
    zI-E_{\rm GW}
  \right)^{-1}
  A^\dagger
  \\
  &=
  \frac{
    AA^\dagger
  }{z}
  +
  \frac{
    AE_{\rm GW}A^\dagger
  }{z^2}
  +
  \mathcal O(z^{-3})
  \end{aligned}
  \label{eq:SigmaGW-high-frequency}
\end{equation}

The dynamic correlation self-energy contains no constant term as
$|z|\rightarrow\infty$.

Consequently, it first contributes to the \emph{second} spectral
moment of the Green's function rather than to the first.


\subsubsection{First moment of the existing MR-$GW$ approximation}

For the practical MR-$GW$ method,
\begin{equation}
  K_B=0
\end{equation}

Equation~\eqref{eq:M1-main-result} therefore becomes
\begin{equation}
  M_1^{\rm MR-GW}
  =
  M_1^D
  +
  \Sigma_{\rm stat}
  \label{eq:M1-existing-MRGW}
\end{equation}

The existing MR-$GW$ approximation consequently omits any
first-moment contribution associated with the mixed blocks
$\Sigma^{12}$ and $\Sigma^{21}$.

This provides an alternative characterization of the information
discarded when these blocks are set to zero.


\subsubsection{High-frequency structure of the mixed blocks}

The retarded screened mixed block has the spectral representation
\begin{equation}
  [\Sigma_{\rm scr}^{12,R}]_{Px}(z)
  =
  \sum_\alpha
  \frac{
    B_{P\alpha}^{\rm scr}
    [d_\alpha^\dagger]_x
  }{
    z-\kappa_\alpha
  }
\end{equation}

At large $|z|$,
\begin{equation}
  [\Sigma_{\rm scr}^{12,R}]_{Px}(z)
  =
  \frac{1}{z}
  \sum_\alpha
  B_{P\alpha}^{\rm scr}
  [d_\alpha^\dagger]_x
  +
  \mathcal O(z^{-2})
\end{equation}
and similarly for $\Sigma^{21}$.

Since
\begin{equation}
  G_D(z)
  =
  \frac{I}{z}
  +
  \mathcal O(z^{-2})
\end{equation}
the contributions
\begin{equation}
  G_D\Sigma^{12}_{\rm scr}
  \qquad\text{and}\qquad
  \Sigma^{21}_{\rm scr}G_D
\end{equation}
begin at order
\begin{equation}
  \frac{1}{z^2}
\end{equation}

They therefore contribute directly to the first spectral moment.

Using the Hermitian representation derived in Sec.~\ref{sec:hermitian-completion}, the sum of these
contributions is exactly
\begin{equation}
  \Delta M_1^{\rm mix}
  =
  T_DK_BT_D^\dagger
  \label{eq:M1-mixed-contribution}
\end{equation}


\subsubsection{Perturbative expansion around the Dyall Hamiltonian}

Introduce again the bookkeeping parameter
\begin{equation}
  \hat H(\lambda)
  =
  \hat H_D
  +
  \lambda\hat V_R
\end{equation}

The exact first moment has the expansion
\begin{equation}
  M_1^{\rm exact}(\lambda)
  =
  M_1^D
  +
  \lambda M_1^{(1)}
  +
  \lambda^2 M_1^{(2)}
  +
  \cdots
  \label{eq:M1-exact-lambda}
\end{equation}

From Sec.~\ref{sec:first-order-blocks}, the exact first-order Green's function is
\begin{equation}
  G_1
  =
  G_D\Sigma_1^{11}G_D
  +
  G_D\Sigma_1^{12}
  +
  \Sigma_1^{21}G_D
\end{equation}
where
\begin{equation}
  \Sigma_1^{22}=0
\end{equation}
for the Dyall Hamiltonian.

Taking the coefficient of $1/z^2$ gives
\begin{equation}
  M_1^{(1)}
  =
  \Sigma_1^{11}
  +
  T_DK_B^{(1)}T_D^\dagger
  \label{eq:M1-firstorder-exact}
\end{equation}
where $K_B^{(1)}$ is constructed from the bare mixed residues,
\begin{equation}
  B^{(1)}
  =
  \bar v_R\mathcal C_D
\end{equation}

Thus the first-order contribution from the mixed Hall blocks is
required by the exact first-moment sum rule.


\subsubsection{First-order exactness of the Hermitian extension}

For the screened kernel,
\begin{equation}
  \mathcal I_D^{\rm stat}
  =
  v_R-[W_D(0)]^x
\end{equation}
the screened interaction has the residual-coupling expansion
\begin{equation}
  W_D(0)
  =
  \lambda v_R
  +
  \lambda^2
  v_RP_D(0)v_R
  +
  \mathcal O(\lambda^3)
  \label{eq:WD-lambda-expansion}
\end{equation}

Hence
\begin{equation}
  \mathcal I_D^{\rm stat}
  =
  \lambda\bar v_R
  +
  \mathcal O(\lambda^2)
\end{equation}
and consequently
\begin{equation}
  K_B
  =
  \lambda K_B^{(1)}
  +
  \mathcal O(\lambda^2)
  \label{eq:KB-firstorder}
\end{equation}

Using Eq.~\eqref{eq:M1-main-result},
\begin{equation}
  M_1^{\rm mix-MR-GW}
  =
  M_1^D
  +
  \lambda
  \left[
    \Sigma_1^{11}
    +
    T_DK_B^{(1)}T_D^\dagger
  \right]
  +
  \mathcal O(\lambda^2)
\end{equation}

Comparison with Eq.~\eqref{eq:M1-firstorder-exact} gives
\begin{equation}
  M_1^{\rm mix-MR-GW}
  =
  M_1^{\rm exact}
  +
  \mathcal O(\lambda^2)
  \label{eq:M1-firstorder-exact-method}
\end{equation}

The proposed Hermitian screened-mixed extension is therefore
\emph{exact through first order in the residual interaction} with
respect to the first spectral moment.


\subsubsection{Why the existing MR-$GW$ method is not first-order exact for
the first moment}

Using Eq.~\eqref{eq:M1-existing-MRGW},
\begin{equation}
  M_1^{\rm MR-GW}
  =
  M_1^D
  +
  \lambda\Sigma_1^{11}
  +
  \mathcal O(\lambda^2)
\end{equation}

The exact result is instead
\begin{equation}
  M_1^{\rm exact}
  =
  M_1^D
  +
  \lambda
  \left[
    \Sigma_1^{11}
    +
    T_DK_B^{(1)}T_D^\dagger
  \right]
  +
  \mathcal O(\lambda^2)
\end{equation}

Therefore
\begin{equation}
  M_1^{\rm MR-GW}
  -
  M_1^{\rm exact}
  =
  -
  \lambda
  T_DK_B^{(1)}T_D^\dagger
  +
  \mathcal O(\lambda^2)
  \label{eq:M1-MRGW-error}
\end{equation}

Unless the mixed contribution happens to vanish for a particular
system or matrix element, practical MR-$GW$ therefore violates the
exact first-moment sum rule already at first residual order.

Restoring the mixed sector removes precisely this leading error.


\subsubsection{Higher-order content of the screened mixed moment}

The static screened kernel contains higher-order terms,
\begin{equation}
  \mathcal I_D^{\rm stat}
  =
  \lambda\bar v_R
  -
  \lambda^2
  [v_RP_D(0)v_R]^x
  +
  \cdots
\end{equation}

Consequently,
\begin{equation}
  T_DK_BT_D^\dagger
\end{equation}
contains a selected sequence of higher-order contributions to the
first moment.

These contributions should \emph{not} be interpreted as the complete
exact second- or higher-order first moment.

In particular, the exact first moment depends on the interacting
one-body density matrix,
\begin{equation}
  \gamma^{\rm exact}
  =
  \gamma_D
  +
  \lambda\gamma^{(1)}
  +
  \lambda^2\gamma^{(2)}
  +
  \cdots
\end{equation}
whereas the one-shot MR-$GW$ construction does not determine all such
density-response terms self-consistently.

Therefore,
\begin{equation}
  \text{the screened Hermitian approximation is first-order exact in
  }M_1
  \text{ but not generally exact beyond first order.}
  \label{eq:M1-order-statement}
\end{equation}


\subsubsection{The role of the dynamic MR-$GW$ bath in higher moments}

Although the MR-$GW$ bath does not enter $M_1$, it enters the next
moment.

From
\begin{equation}
  M_2
  =
  \mathbb T
  \mathbb H^2
  \mathbb T^\dagger
\end{equation}
one obtains
\begin{equation}
  M_2
  =
  T_DK_{\rm top}^2T_D^\dagger
  +
  AA^\dagger
  \label{eq:M2-preview}
\end{equation}

Thus the dynamic MR-$GW$ amplitudes first appear through
\begin{equation}
  AA^\dagger
\end{equation}
in the second spectral moment.

This provides a useful interpretation of the division of roles:
\begin{equation}
  \begin{array}{rcl}
  M_0
  &:&
  \text{normalization}
  \\[1mm]
  M_1
  &:&
  \text{static and mixed-reference couplings}
  \\[1mm]
  M_2\text{ and higher}
  &:&
  \text{dynamic MR-}GW\text{ screening enters explicitly}
  \end{array}
  \label{eq:moment-role-summary}
\end{equation}


\subsubsection{First-moment error as a practical diagnostic}

Equation~\eqref{eq:exact-first-moment-commutator} suggests a useful
numerical diagnostic.

Given an approximate one-body density matrix $\gamma$ obtained from the
resulting Green's function, one may independently construct
\begin{equation}
  M_1^{\rm comm}[\gamma]
\end{equation}
from the equal-time commutator with the physical Hamiltonian.

This can be compared with the spectral moment
\begin{equation}
  M_1^{\rm spec}
  =
  \int
  d\omega\
  \omega A(\omega)
\end{equation}

The difference
\begin{equation}
  \Delta_1
  =
  M_1^{\rm spec}
  -
  M_1^{\rm comm}[\gamma]
  \label{eq:first-moment-diagnostic}
\end{equation}
provides a compact measure of the internal moment consistency of a
given truncation.

For the exact Green's function,
\begin{equation}
  \Delta_1=0
\end{equation}

This comparison may be particularly useful when assessing the three
candidate mixed kernels
\begin{equation}
  \bar v_R
  \qquad
  v_R-W_D^x(0)
  \qquad
  W_D(0)-W_D^x(0)
\end{equation}
%%%%%%%%%%%%%%%%%%%%%%%%%
\subsection{Computational formulation and scaling}
\label{sec:implementation}
%%%%%%%%%%%%%%%%%%%%%%%%%

\Proposed This subsection turns the formal construction into a practical
algorithm.  Costs quoted for the parent spectral MR-$GW$ implementation are
those of Ref.~\cite{WangFangLi2026}; the additional mixed-block costs are
estimates for the implementation proposed here.


We finally discuss the practical implementation of the screened-mixed
MR-$GW$ approximation developed in the preceding sections.

The objective is not to design a fully optimized low-scaling
implementation at this stage, but rather to establish that the new
mixed-block contribution can be added to the existing MR-$GW$
algorithm at moderate additional cost.

We follow the notation of Ref.~\cite{WangFangLi2026}.  Let
\begin{equation}
  N_{\rm orb}
  =
  N_c+N_a+N_v
\end{equation}
where $N_c$, $N_a$, and $N_v$ denote the number of core, active, and
virtual spin orbitals, respectively.

We further define
\begin{equation}
  N_I
  =
  N_c+N_v
\end{equation}
as the number of inactive orbitals.

Let
\begin{equation}
  D_0
  \qquad
  D_+
  \qquad
  D_-
\end{equation}
be the determinant dimensions of the neutral, electron-attached, and
electron-removed active sectors, and define
\begin{equation}
  D
  =
  \max(D_0,D_+,D_-)
\end{equation}

The dimension of each MR-RPA $A$ or $B$ block is denoted by $K$, with
the bound
\begin{equation}
  K
  \le
  N_cN_v
  +
  D(N_c+N_v+1)
  \label{eq:K-dimension}
\end{equation}

Finally, let $L$ denote the number of charged channels retained in the
Dyall one-particle Green's function.  For a full spectral representation,
\begin{equation}
  L
  \le
  N_c+N_v+D_++D_-
  \le
  N_c+N_v+2D
  \label{eq:L-dimension}
\end{equation}
The second inequality is the bound used in Ref.~\cite{WangFangLi2026}; the
first allows for symmetry-forbidden or otherwise omitted charged channels.


\subsubsection{Quantities already required by MR-$GW$}

The existing spectral MR-$GW$ implementation already provides the
following quantities:
\begin{equation}
  E_{0,N_a}^A
  \qquad
  E_{\mu,N_a}^A
  \qquad
  E_{\mu,N_a+1}^A
  \qquad
  E_{\mu,N_a-1}^A
\end{equation}
together with the transition amplitudes
\begin{subequations}
  \begin{align}
    D_{\mu,x}^{[+1]}
    &=
    \braket{
      \Xi_{\mu,N_a+1}^A
      |
      \hat x^\dagger
      |
      \Xi_{0,N_a}^A
    }
    \\
    D_{\mu,x}^{[-1]}
    &=
    \braket{
      \Xi_{\mu,N_a-1}^A
      |
      \hat x
      |
      \Xi_{0,N_a}^A
    }
    \\
    D_{\mu,xy}^{[0]}
    &=
    \braket{
      \Xi_{\mu,N_a}^A
      |
      \hat x^\dagger\hat y
      |
      \Xi_{0,N_a}^A
    }
  \end{align}
\end{subequations}

These quantities are used to construct the MR-RPA matrices, whose
diagonalization yields
\begin{equation}
  \Omega_I
  \qquad
  X_{\mu I}
  \qquad
  Y_{\mu I}
\end{equation}

The screened-interaction amplitudes are
\begin{equation}
  M_{pr,I}
  =
  v^R_{pr,qs}
  \left[
    \braket{
      \Phi_0
      |
      \hat q^\dagger\hat s
      |
      \Phi_\mu
    }
    X_{\mu I}
    +
    \braket{
      \Phi_\mu
      |
      \hat q^\dagger\hat s
      |
      \Phi_0
    }
    Y_{\mu I}
  \right]
  \label{eq:M-implementation}
\end{equation}

No modification of these parts of the MR-$GW$ calculation is required
by the present extension.


\subsubsection{Cost of the parent spectral MR-$GW$ calculation}

For reference, the dense spectral implementation of
Ref.~\cite{WangFangLi2026} has the following principal costs:
\begin{center}
\begin{tabular}{l c}
\hline
Operation & Leading cost \\
\hline
Four-index AO--MO transformation
&
$\mathcal O(N_{\rm orb}^5)$
\\
Active-sector diagonalizations
&
$\mathcal O(D^3)$
\\
Active transition densities
&
$\mathcal O(D^2N_a^2)$
\\
Dense MR-RPA diagonalization
&
$\mathcal O(K^3)$
\\
Construction of all $M_{pr,I}$
&
$\mathcal O(N_{\rm orb}^2K^2)$
\\
MR-$GW$ self-energy on a frequency grid
&
$\mathcal O(N_\omega N_\Sigma K L)$
\\
Physical Dyson inversion
&
$\mathcal O(N_\omega N_{\rm orb}^3)$
\\
\hline
\end{tabular}
\end{center}

Without an active space,
\begin{equation}
  K
  =
  \mathcal O(N_{\rm orb}^2)
\end{equation}
so the dense spectral implementation has the conventional
\begin{equation}
  \mathcal O(N_{\rm orb}^6)
\end{equation}
scaling of a full-pole dense $G_0W_0$ implementation.

At fixed active space, the same asymptotic orbital-space scaling is
retained in MR-$GW$, with additional determinant-space costs associated
with the active reference.


\subsubsection{Static screened interaction required by the new method}

The new mixed kernel requires only the zero-frequency screened
interaction,
\begin{equation}
  W_D^0
  =
  W_D(0)
\end{equation}

Using the MR-RPA spectral representation,
\begin{equation}
  [w_D^c]_{pr,qs}(\omega)
  =
  \sum_{I>0}
  \left[
    \frac{
      M_{pr,I}M_{qs,I}^*
    }{
      \omega-\Omega_I+i0^+
    }
    -
    \frac{
      M_{pr,I}^*M_{qs,I}
    }{
      \omega+\Omega_I-i0^+
    }
  \right]
\end{equation}
one obtains at zero frequency
\begin{equation}
  [w_D^c(0)]_{pr,qs}
  =
  -
  \sum_{I>0}
  \frac{
    M_{pr,I}M_{qs,I}^*
    +
    M_{pr,I}^*M_{qs,I}
  }{
    \Omega_I
  }
  \label{eq:wc-zero}
\end{equation}

For real quantities,
\begin{equation}
  [w_D^c(0)]_{pr,qs}
  =
  -2
  \sum_{I>0}
  \frac{
    M_{pr,I}M_{qs,I}
  }{
    \Omega_I
  }
  \label{eq:wc-zero-real}
\end{equation}

Only the mixed slice
\begin{equation}
  [W_D^0]_{Pw,yz}
\end{equation}
is required.

If this slice is formed explicitly from the already available
$M_{pr,I}$, its cost is
\begin{equation}
  \mathcal O
  \left(
    K N_I N_a^3
  \right)
  \label{eq:Wmixed-cost}
\end{equation}
with storage
\begin{equation}
  \mathcal O
  \left(
    N_I N_a^3
  \right)
  \label{eq:Wmixed-memory}
\end{equation}

At fixed active space this is only cubic in the external orbital
dimension and is therefore subleading relative to the dense MR-RPA
step.


\subsubsection{Avoiding the explicit three-operator transition tensor}

A direct implementation of the screened mixed residues might first
construct
\begin{equation}
  \mathcal C_{\mu;z,yw}^{D,+}
\end{equation}
and
\begin{equation}
  \mathcal C_{\mu;z,yw}^{D,-}
\end{equation}

Their storage would scale as
\begin{equation}
  \mathcal O(DN_a^3)
\end{equation}
and transformation to all dense charged eigenstates would require at
most
\begin{equation}
  \mathcal O(D^2N_a^3)
\end{equation}
work.

This explicit tensor is unnecessary.

Define instead, for every inactive orbital $P$, the active-space
operator
\begin{equation}
  \hat Q_P
  =
  \sum_{zyw}
  [\mathcal I_D^{\rm stat}]_{Pz,yw}
  \left[
    \frac12
    \hat y^\dagger\hat w\hat z
    -
    \gamma_{yw}^D
    \hat z
  \right]
  \label{eq:QP-operator}
\end{equation}

The screened mixed residues can then be evaluated directly as
\begin{subequations}
  \label{eq:B-source-form}
  \begin{align}
    B_{P\mu}^{+,{\rm scr}}
    &=
    \braket{
      \Xi_{0,N_a}^A
      |
      \hat Q_P
      |
      \Xi_{\mu,N_a+1}^A
    }
    \\
    B_{P\mu}^{-,{\rm scr}}
    &=
    \braket{
      \Xi_{\mu,N_a-1}^A
      |
      \hat Q_P
      |
      \Xi_{0,N_a}^A
    }
  \end{align}
\end{subequations}

Equivalently, introduce the charged source vectors
\begin{subequations}
  \begin{align}
    \ket{\chi_P^+}
    &=
    \hat Q_P^\dagger
    \ket{\Xi_{0,N_a}^A}
    \\
    \ket{\chi_P^-}
    &=
    \hat Q_P
    \ket{\Xi_{0,N_a}^A}
  \end{align}
\end{subequations}

Then
\begin{subequations}
  \begin{align}
    B_{P\mu}^{+,{\rm scr}}
    &=
    \braket{
      \chi_P^+
      |
      \Xi_{\mu,N_a+1}^A
    }
    \\
    B_{P\mu}^{-,{\rm scr}}
    &=
    \braket{
      \Xi_{\mu,N_a-1}^A
      |
      \chi_P^-
    }
  \end{align}
\end{subequations}

This formulation avoids storing any object with indices
$(\mu,z,y,w)$.


\subsubsection{Cost of the contracted-source construction}

A general three-active-index operator connects a determinant to at most
$\mathcal O(N_a^3)$ determinants.

Constructing one source vector therefore requires at most
\begin{equation}
  \mathcal O(DN_a^3)
\end{equation}
operations.

For all inactive orbitals,
\begin{equation}
  \mathcal O(N_I DN_a^3)
\end{equation}
work is sufficient to form the source vectors.

Projection of each source vector onto all charged eigenstates requires
at most
\begin{equation}
  \mathcal O(D^2)
\end{equation}
work per inactive orbital.

Hence a conservative upper bound for direct construction of all mixed
residues is
\begin{equation}
  \mathcal O
  \left[
    N_I
    \left(
      DN_a^3+D^2
    \right)
  \right]
  \label{eq:B-direct-cost}
\end{equation}

The resulting residues require only
\begin{equation}
  \mathcal O
  \left[
    N_I(D_++D_-)
  \right]
  \label{eq:B-memory}
\end{equation}
storage.

If only a subset of charged states is required, $D$ in the projection
step can be replaced by the number of retained charged eigenstates.


\subsubsection{Factorized evaluation of the screened contribution}

The correlation part of the static screened interaction has a low-rank
MR-RPA representation.

For the crossed slice entering the mixed kernel,
\begin{equation}
  [w_D^c(0)]_{Pw,yz}
  =
  -
  2
  \sum_{I>0}
  \frac{
    M_{Pw,I}M_{yz,I}
  }{
    \Omega_I
  }
\end{equation}
for real quantities.

Consequently, it is not necessary to construct
$[W_D^0]_{Pw,yz}$ explicitly.

The contraction with the connected active tensor may instead be
factorized by first defining
\begin{equation}
  F_{\mu w I}^{\sigma}
  =
  \sum_{yz}
  M_{yz,I}
  \mathcal C_{\mu;z,yw}^{D,\sigma}
  \label{eq:F-factor}
\end{equation}

The correlation part of the residue then takes the form
\begin{equation}
  B_{P\mu}^{\sigma,c}
  =
  2
  \sum_{Iw}
  \frac{
    M_{Pw,I}
    F_{\mu w I}^{\sigma}
  }{
    \Omega_I
  }
  \label{eq:B-factorized}
\end{equation}
up to the sign associated with the kernel convention of
Sec.~\ref{sec:static-mixed-kernel}.

A direct factorized implementation therefore has the schematic cost
\begin{equation}
  \mathcal O
  \left(
    KDN_a^3
    +
    N_IKDN_a
  \right)
  \label{eq:B-factorized-cost}
\end{equation}
plus the charged-state projection cost when the transition tensors are
generated from source vectors.

The MR-RPA mode index $I$ can be blocked, so the intermediate
$F_{\mu w I}^{\sigma}$ need never be stored for all modes
simultaneously.


\subsubsection{Construction and storage of \texorpdfstring{$K_B$}{KB}}

Once the screened residues have been obtained, the Hermitian mixed
matrix is defined by
\begin{equation}
  [K_B]_{P\alpha}
  =
  B_{P\alpha}^{\rm scr}
  \qquad
  [K_B]_{\alpha P}
  =
  B_{P\alpha}^{{\rm scr}*}
\end{equation}

The only nonzero elements connect inactive poles to active charged
poles.

Therefore,
\begin{equation}
  N_{\rm nz}(K_B)
  =
  \mathcal O
  \left[
    N_I(D_++D_-)
  \right]
  \label{eq:KB-nnz}
\end{equation}
and the matrix-vector product
\begin{equation}
  K_Bx
\end{equation}
requires the same order of work.

Thus $K_B$ is a very sparse object even though the complete auxiliary
Hamiltonian is large.


\subsubsection{Size of the MR-$GW$ pole bath}

The spectral MR-$GW$ self-energy contains, for every positive MR-RPA
mode $I$, the charged channels
\begin{equation}
  a
  \qquad
  i
  \qquad
  (\mu,+)
  \qquad
  (\mu,-)
\end{equation}

The number of self-energy poles is therefore bounded by
\begin{equation}
  N_{\rm bath}
  \lesssim
  K L
  \label{eq:Nbath}
\end{equation}

The formal Hermitian matrix of Sec.~\ref{sec:hermitian-completion} consequently has dimension
\begin{equation}
  N_H
  =
  L+N_{\rm bath}
  =
  \mathcal O(KL)
  \label{eq:Hdimension}
\end{equation}

A dense construction or dense diagonalization of this matrix would be
prohibitively expensive and is neither necessary nor recommended.


\subsubsection{Matrix-free application of the common Hermitian matrix}

The common Hermitian matrix is
\begin{equation}
  \mathbb H
  =
  \begin{pmatrix}
    K_D
    +
    T_D^\dagger
    \Sigma_{\rm stat}
    T_D
    +
    K_B
    &
    T_D^\dagger A
    \\
    A^\dagger T_D
    &
    E_{\rm GW}
  \end{pmatrix}
\end{equation}

For a trial vector
\begin{equation}
  x
  =
  \begin{pmatrix}
    x_D\\
    x_{\rm bath}
  \end{pmatrix}
\end{equation}
its action may be evaluated as
\begin{subequations}
  \begin{align}
    y_D
    ={}&
    K_Dx_D
    +
    K_Bx_D
    +
    T_D^\dagger
    \Sigma_{\rm stat}
    (T_Dx_D)
    +
    T_D^\dagger
    (Ax_{\rm bath})
    \\
    y_{\rm bath}
    ={}&
    E_{\rm GW}x_{\rm bath}
    +
    A^\dagger(T_Dx_D)
  \end{align}
\end{subequations}

Thus neither
\begin{equation}
  T_D^\dagger A
\end{equation}
nor the full matrix $\mathbb H$ needs to be stored.

If the self-energy amplitudes $A$ are treated as an explicit
rectangular matrix, a matrix-vector product requires approximately
\begin{equation}
  \mathcal O
  \left(
    N_{\rm orb} K L
  \right)
  \label{eq:H-matvec-cost}
\end{equation}
work, in addition to the much cheaper sparse $K_B$ multiplication.

At fixed active space,
\begin{equation}
  K=\mathcal O(N_{\rm orb}^2)
  \qquad
  L=\mathcal O(N_{\rm orb})
\end{equation}
so
\begin{equation}
  C_{\rm matvec}
  =
  \mathcal O(N_{\rm orb}^4)
\end{equation}

This remains below the
$\mathcal O(N_{\rm orb}^6)$ cost of the dense MR-RPA spectral
implementation.


\subsubsection{Factorized application of the MR-$GW$ bath}

It is not necessary to form $A$ explicitly either.

The entries of $A$ factor into the already available screened
interaction amplitudes $M_{pr,I}$ and the Dyall addition/removal
transition densities.

For example, the active electron-addition amplitudes contain
\begin{equation}
  D_{\mu,x}^{[+1]}M_{xp,I}
\end{equation}
while the removal amplitudes contain
\begin{equation}
  D_{\mu,x}^{[-1]}M_{px,I}
\end{equation}

Therefore the actions
\begin{equation}
  Ax_{\rm bath}
\end{equation}
and
\begin{equation}
  A^\dagger y
\end{equation}
can be implemented directly from the same precontracted objects used
to evaluate the MR-$GW$ self-energy.

The Hermitian formulation therefore does not require storing an
additional object of dimension
\begin{equation}
  N_{\rm orb}\times KL
\end{equation}


\subsubsection{Two practical solution strategies}

The common Hermitian representation can be exploited in two
complementary ways.

For selected charged excitations, one may apply Lanczos, Davidson, or
another iterative Hermitian eigensolver directly to $\mathbb H$ using
only the matrix-vector products described above.

For a frequency-dependent spectral function, it is preferable to
eliminate the MR-$GW$ bath analytically.  The Schur complement gives
\begin{equation}
  \begin{aligned}
  R_{\rm top}^{-1}(z)
  ={}&
  zI
  -
  K_D
  -
  T_D^\dagger\Sigma_{\rm stat}T_D
  -
  K_B
  \\
  &
  -
  T_D^\dagger
  A
  (zI-E_{\rm GW})^{-1}
  A^\dagger
  T_D
  \end{aligned}
  \label{eq:top-schur-practical}
\end{equation}

The physical Green's function is then
\begin{equation}
  G^R(z)
  =
  T_D
  R_{\rm top}(z)
  T_D^\dagger
  \label{eq:G-top-practical}
\end{equation}

Equation~\eqref{eq:top-schur-practical} involves only an
$L$-dimensional top space.  The large bath enters exclusively through
the action of the factorized dynamic self-energy.

This is likely the preferred route for computing spectral functions on
a frequency grid.


\subsubsection{Scaling of the new ingredients}

The principal additional operations introduced relative to the parent
MR-$GW$ calculation are summarized below:
\begin{center}
\begin{tabular}{l c c}
\hline
New operation
&
Work
&
Storage
\\
\hline
Explicit mixed $W_D(0)$ slice
&
$\mathcal O(KN_IN_a^3)$
&
$\mathcal O(N_IN_a^3)$
\\
Explicit charged three-operator tensors
&
$\mathcal O(D^2N_a^3)$
&
$\mathcal O(DN_a^3)$
\\
Contracted charged source vectors
&
$\mathcal O[N_I(DN_a^3+D^2)]$
&
$\mathcal O(N_ID)$
\\
Factorized screened contraction
&
$\mathcal O(KDN_a^3+N_IKDN_a)$
&
blockable
\\
Construction/storage of $K_B$
&
$\mathcal O(N_ID)$
&
$\mathcal O(N_ID)$
\\
Hermitian matrix-vector product
&
$\mathcal O(N_{\rm orb}KL)$
&
matrix-free
\\
\hline
\end{tabular}
\end{center}

The explicit three-operator tensor and the explicit $W_D(0)$ slice are
alternative implementation routes rather than mandatory simultaneous
costs.

The recommended implementation uses contracted source vectors and
factorized MR-RPA amplitudes wherever possible.


\subsubsection{Asymptotic orbital scaling at fixed active space}

For fixed $N_a$ and fixed determinant dimensions,
\begin{equation}
  D=\mathcal O(1)
\end{equation}
while
\begin{equation}
  N_I
  =
  \mathcal O(N_{\rm orb})
\end{equation}
\begin{equation}
  L
  =
  \mathcal O(N_{\rm orb})
\end{equation}
and
\begin{equation}
  K
  =
  \mathcal O(N_{\rm orb}^2)
\end{equation}

The new static mixed-screening operations then scale no worse than
approximately
\begin{equation}
  \mathcal O(N_{\rm orb}^3)
\end{equation}
while a matrix-free application of the final Hermitian problem scales
as
\begin{equation}
  \mathcal O(N_{\rm orb}^4)
\end{equation}

Both are subleading relative to the
\begin{equation}
  \mathcal O(N_{\rm orb}^6)
\end{equation}
dense MR-RPA diagonalization already present in the spectral
implementation.

Consequently,
\begin{equation}
  \text{the screened-mixed PSD extension does not increase the
  leading orbital-space scaling of dense spectral MR-}GW
  \label{eq:no-scaling-increase}
\end{equation}


\subsubsection{Determinant-space scaling}

The principal genuinely new determinant-space quantity is the
three-active-operator source vector.

Using the contracted-source formulation, its conservative cost is
\begin{equation}
  \mathcal O
  \left[
    N_I(DN_a^3+D^2)
  \right]
\end{equation}

This should be compared with the existing active-sector diagonalization
cost
\begin{equation}
  \mathcal O(D^3)
\end{equation}
and MR-RPA determinant dependence through
\begin{equation}
  K
  \le
  N_cN_v
  +
  D(N_c+N_v+1)
\end{equation}

For the small and moderate complete active spaces targeted by the
initial implementation, this additional cost should therefore remain
manageable.

For substantially larger active spaces, the same source-vector
formulation is naturally compatible with iterative CI or DMRG
representations because it only requires the action of
$\hat Q_P$ on the reference wave function and overlaps with charged
response states.


\subsubsection{Relation to a future frequency-domain implementation}

The spectral implementation is useful for developing and validating
the present theory, but it should not be regarded as the ultimate
algorithm.

Ref.~\cite{WangFangLi2026} already describes an alternative
frequency-domain formulation in which the active neutral and charged
responses are obtained from sparse linear equations rather than full
diagonalization.

In that formulation the active response scales schematically as
\begin{equation}
  \mathcal O
  \left(
    N_\omega
    N_{\rm it}
    D
    N_a^5
    N_{\rm orb}
  \right)
\end{equation}
while modern RI-based orbital-space contractions can reduce selected
MR-$GW$ quantities to the same asymptotic scaling as modern
single-reference $GW$.

The mixed source vectors introduced here are particularly compatible
with such an approach because they are already expressed as operator
actions,
\begin{equation}
  \hat Q_P
  \ket{\Xi_0}
\end{equation}
rather than requiring complete charged-state eigenvectors.

Thus a future low-scaling formulation can evaluate the screened mixed
sector directly through charged-sector resolvents without constructing
all $B_{P\mu}^{\sigma}$ explicitly.


\subsubsection{Recommended initial implementation}

For a first proof-of-principle implementation, the following workflow
is sufficient:
\begin{enumerate}
  \item Perform CASCI or CASSCF and construct the Dyall Hamiltonian.

  \item Diagonalize the neutral and $N_a\pm1$ active sectors, as in the
  existing MR-$GW$ implementation.

  \item Construct the standard transition densities
  $D^{[0]}$, $D^{[+1]}$, and $D^{[-1]}$.

  \item Solve the dense MR-RPA problem and obtain
  $\Omega_I$, $X_I$, $Y_I$, and $M_{pr,I}$.

  \item Construct only the static mixed kernel required for
  \begin{equation}
    \mathcal I_D^{\rm stat}
    =
    v_R-W_D^x(0)
  \end{equation}

  \item Build the contracted source operators $\hat Q_P$ and obtain
  the screened residues $B_{P\mu}^{\pm,\rm scr}$ without forming the
  full three-operator transition tensor.

  \item Assemble the sparse Hermitian matrix $K_B$.

  \item Reuse the existing MR-$GW$ pole amplitudes to apply the dynamic
  bath matrix-free.

  \item Compute either selected poles with an iterative Hermitian
  eigensolver or the spectral function through the top-space Schur
  complement of Eq.~\eqref{eq:top-schur-practical}.

  \item Verify numerically the PSD property, zeroth spectral moment,
  first spectral moment, and all limiting cases derived in the
  preceding sections.
\end{enumerate}


\subsubsection{Quantities that should not be formed explicitly}

Although the formal derivation uses several large matrices, an
implementation should avoid explicitly constructing
\begin{equation}
  \mathbb H_{\rm mix-MR-GW}
\end{equation}
\begin{equation}
  T_D^\dagger A
\end{equation}
the complete $N_{\rm orb}\times KL$ self-energy amplitude matrix, and
the complete three-operator transition tensor
\begin{equation}
  \mathcal C_{\mu;z,yw}^{D,\sigma}
\end{equation}

All of these objects admit factorized or matrix-free actions.

The only genuinely new object that should normally be stored is the
sparse inactive--active coupling matrix
\begin{equation}
  K_B
\end{equation}
%%%%%%%%%%%%%%%%%%%%%%%%%
\section{Consistency limits and hierarchy of approximations}
\label{sec:consistency}
%%%%%%%%%%%%%%%%%%%%%%%%%

The formal hierarchy and the proposed approximation are constrained by a set
of limiting cases that should be regarded as non-negotiable implementation
tests.
%%%%%%%%%%%%%%%%%%%%%%%%%
\subsection{No residual interaction and full active space}
%%%%%%%%%%%%%%%%%%%%%%%%%

\Exact If
\begin{equation}
  \hat V_R=0
\end{equation}
then all residual self-energy and screening corrections vanish and
\begin{equation}
  G=G_D
\end{equation}
The Dyall cumulants may remain nonzero and continue to enter response
functions, but they do not generate an additional physical two-point
self-energy in the absence of a residual interaction.

If the active space spans the complete one-particle space, then
\begin{equation}
  \hat H_D=\hat H
  \qquad
  \hat V_R=0
\end{equation}
so that
\begin{equation}
  G_D=G_{\rm exact}
\end{equation}
%%%%%%%%%%%%%%%%%%%%%%%%%
\subsection{Gaussian or empty-active-space limit}
%%%%%%%%%%%%%%%%%%%%%%%%%

\Exact When the Dyall reference becomes Gaussian,
\begin{equation}
  C_{n\ge2}^D=0
\end{equation}
The auxiliary cumulant hierarchy disappears,
\begin{equation}
  \Sigma^{12}
  =
  \Sigma^{21}
  =
  \Sigma^{22}
  =
  0
\end{equation}
while
\begin{equation}
  \Gamma_D\to1
  \qquad
  P_D\to-iG_DG_D
\end{equation}
The generalized Dyson equation therefore reduces to the ordinary Dyson
equation and the formal hierarchy collapses to conventional Hedin theory.
For an empty active space, the practical MR-$GW$ construction correspondingly
reduces to the single-reference $G_0W_0$ limit associated with the chosen
one-particle reference.
%%%%%%%%%%%%%%%%%%%%%%%%%
\subsection{Recovery of the published MR-$GW$ scheme}
%%%%%%%%%%%%%%%%%%%%%%%%%

\Established In the proposed Hermitian representation, setting
\begin{equation}
  K_B=0
\end{equation}
removes the mixed sector exactly and recovers the Wang--Fang--Li MR-$GW$
Dyson equation and its pole-space linearization.
%%%%%%%%%%%%%%%%%%%%%%%%%
\subsection{Bare mixed-block limit}
%%%%%%%%%%%%%%%%%%%%%%%%%

\Exact At leading residual order,
\begin{equation}
  W_D(0)\to v_R
\end{equation}
implies
\begin{equation}
  \mathcal I_D^{\rm stat}
  \to
  \bar v_R
\end{equation}
The screened residues therefore reduce to the exact first-order mixed
self-energy blocks derived in Sec.~\ref{sec:first-order-blocks}.  Since
$\Sigma_1^{22}=0$ for the Dyall partition, the Hermitian mixed extension
recovers the complete first-order Green's function in this limit.
%%%%%%%%%%%%%%%%%%%%%%%%%
\subsection{PSD, causality, and spectral moments}
%%%%%%%%%%%%%%%%%%%%%%%%%

\Proposed For the finite Hermitian completion of
Sec.~\ref{sec:hermitian-completion}, Hermiticity guarantees a retarded causal
resolvent with residue matrices of the form
\begin{equation}
  Z_\lambda Z_\lambda^\dagger\succeq0
\end{equation}
Consequently,
\begin{equation}
  A(\omega)\succeq0
  \qquad
  \int d\omega\,A(\omega)=I
\end{equation}
The first moment satisfies
\begin{equation}
  M_1
  =
  M_1^D
  +
  \Sigma_1^{11}
  +
  T_DK_BT_D^\dagger
\end{equation}
and is exact through first order in the residual interaction.  Higher-order
terms generated by static screening and repeated insertions of $K_B$ define
the chosen approximation and are not claimed to constitute the complete
higher-order Dyall-Hedin expansion.
%%%%%%%%%%%%%%%%%%%%%%%%%
\subsection{Hierarchy}
%%%%%%%%%%%%%%%%%%%%%%%%%

The resulting sequence of approximations can be summarized as
\begin{equation}
  \begin{array}{c}
  \text{exact Dyall-Hedin hierarchy}
  \\[1mm]
  \Downarrow
  \\
  \text{published MR-}GW
  \quad
  \left(
    \Sigma^{12}
    =
    \Sigma^{21}
    =
    \Sigma^{22}
    =
    0
  \right)
  \\[2mm]
  \Downarrow
  \\
  \text{static screened mixed extension}
  \quad
  \left[
    \mathcal I_D^{\rm stat}
    =
    v_R-W_D^x(0)
  \right]
  \\[2mm]
  \Downarrow
  \\
  \text{Hermitian/PSD completion}
  \\[2mm]
  \Downarrow
  \\
  \text{future dynamical MR-}GW\Gamma
  \text{ theory}
  \end{array}
\end{equation}
%%%%%%%%%%%%%%%%%%%%%%%%%
\section{Conclusions and outlook}
\label{sec:conclusion}
%%%%%%%%%%%%%%%%%%%%%%%%%

These notes develop two connected pieces of a multireference Green's-function
program based on a CASSCF wave function and the Dyall Hamiltonian.

For Objective~I, the exact non-Gaussian Dyall reference is embedded in an
enlarged Gaussian source space in which an ordinary matrix Dyson equation and
a superspace Bethe--Salpeter equation can be written.  The non-Gaussian
character of the reference reappears as a hierarchy of auxiliary cumulant
vertices.  The physical projection reproduces Hall's generalized Dyson
equation, the exact source derivative gives
\begin{equation}
  \frac{\delta G_D}{\delta U}
  =
  G_DG_D-C_2^{D,ph}
\end{equation}
and the resulting first-order expansion reproduces the four generalized
self-energy blocks of Wang, Fang, and Li.  Their practical MR-$GW$ method is
then identified as a channel-selective truncation: the correlated Dyall
response is retained in the screening channel, the ordinary $GW$ topology is
used in the physical $11$ self-energy, and the mixed and auxiliary Hall blocks
are discarded.

For Objective~II, the leading discarded mixed blocks are restored through the
static kernel
\begin{equation}
  \mathcal I_D^{\rm stat}
  =
  v_R-W_D^x(0)
\end{equation}
obtained from the parent MR-$GW$ self-energy under frozen-$W$ and static
approximations.  The resulting mixed residues are embedded together with the
dynamic MR-$GW$ pole bath in a common Hermitian auxiliary Hamiltonian.  This
construction reproduces the mixed Green-function terms through first order,
is causal and positive semidefinite, preserves the zeroth spectral moment,
and restores the complete first-order contribution to the first moment.
Importantly, the additional operations do not increase the leading dense
orbital-space scaling of the parent spectral implementation at fixed active
space.

Several extensions are natural.  The first is numerical: implement the bare,
parent-consistent static, and fully screened static mixed kernels and compare
their spectral properties, quasiparticle energies, and moment errors on
systems where the neglected Hall blocks are appreciable.  The second is
formal: replace the static kernel by the full two-frequency object
\begin{equation}
  W_D(\nu)\Gamma_D(\omega,\nu)
\end{equation}
and ultimately restore residual vertex corrections
$\delta W_D/\delta G_D$ and the associated mixed Bethe--Salpeter equations.
A third direction is to clarify how the Hermitian completion reorganizes
higher-order contributions involving $\Sigma^{12}$, $\Sigma^{21}$, and
$\Sigma^{22}$ without identifying that completion with any single Hall block.

The practical philosophy is therefore deliberately incremental:
retain the exact CASSCF/Dyall reference, reuse the MR-RPA screening already
available in MR-$GW$, restore the cheapest missing mixed-sector information
first, and increase the vertex content only when required by accuracy.

%%%%%%%%%%%%%%%%%%%%%%%%
\bibliography{parquet-mr-hedin}
%%%%%%%%%%%%%%%%%%%%%%%%

\end{document}

